% ============================
% CAPÍTULO 5: RESULTADOS
% ============================

\phantomsection
\renewcommand{\tablename}{Tabla}
\setstretch{1.5}
\fontsize{14}{16}\selectfont

\chapter{Resultados}
\label{chap:5_resultados}

\makeatletter
\@afterindentfalse\@afterheading
\makeatother

\setcounter{section}{0}
\renewcommand{\thesection}{\thechapter.\arabic{section}}

\fontsize{12}{14}\selectfont
\vspace{\baselineskip}

\phantomsection
\setstretch{1.5}

En el presente capítulo se exponen los resultados obtenidos a partir de la metodología descrita en el Capítulo~\ref{chap:4_metodologia}. El capítulo se organiza en cinco secciones. La primera (§\ref{sec:extraccion}) describe el proceso de extracción terminológica híbrida: los inventarios validados para los subcorpus en inglés y español, el análisis de patrones morfosintácticos, las propiedades estadísticas de terminologicidad y una caracterización de la incrustación terminológica como propiedad estructural de los inventarios. La segunda (§\ref{sec:modelo_tc}) aborda el modelado supervisado a nivel de token y de candidato mediante cinco arquitecturas de clasificación, el análisis de importancia de rasgos, el protocolo de reserva de tipos terminológicos, la curva de saturación terminológica y los experimentos de tamaño y diversidad tipológica de corpus. La tercera (§\ref{sec:evaluacion_solapamiento}) caracteriza el sesgo estadístico en la evaluación de rasgos distribucionales y valida una corrección mediante recomputación por partición, con evaluación externa sobre el banco de datos ACTER. La cuarta (§\ref{sec:tecla}) describe la arquitectura y evaluación de TECLA, la librería que integra computacionalmente todos los componentes del modelo. La quinta (§\ref{sec:glosario}) presenta el glosario bilingüe resultante y las asimetrías terminológicas entre lenguas.

Estos resultados se articulan en el marco de \texttt{TECLA} (\textit{Terminology Extraction and Corpus Linguistics Application}), que implementa el flujo completo de procesamiento terminológico y materializa computacionalmente el modelo teórico propuesto. Más allá de los inventarios y métricas individuales, un hilo conductor atraviesa todos los resultados del capítulo: la \textbf{diversidad contextual} ---el número de contextos discursivos distintos en que aparece una unidad--- emerge sistemáticamente como el predictor primario de terminologicidad, tanto en los rasgos de los clasificadores supervisados como en la correlación con la detectabilidad de términos no vistos durante el entrenamiento ($\rho_s = 0{,}35$ EN; $\rho_s = 0{,}46$ ES; ambos $p < 0{,}001$), frente a la frecuencia bruta, cuya correlación con la detectabilidad no es significativa en el subcorpus inglés ($p = 0{,}14$). Este resultado constituye una corroboración empírica sólida de uno de los principios nucleares de la teoría terminológica contemporánea (§\ref{sec:evaluacion_solapamiento}): los términos se caracterizan no por su abundancia sino por la variedad y especificidad de sus contextos de uso.
% Chapter 4.1
\phantomsection
\setstretch{1.5}
\fontsize{12}{14}\selectfont

% ============================================================
%  SECCIÓN: Extracción terminológica
%  Requiere en el preámbulo:
%    \usepackage{booktabs, longtable, array, algorithm2e}
%    \usepackage[spanish,ruled,vlined]{algorithm2e}
%  Marcadores: % [FIGURA label]
% ============================================================

\section{Extracción terminológica}
\label{sec:extraccion}

La extracción terminológica se aplicó de forma independiente sobre los dos
subcorpus que integran el corpus de trabajo: el subcorpus en inglés, compuesto
por 50~documentos (artículos de revista, capítulos de libro y monografías)
sobre lingüística de corpus en el ámbito anglosajón, y el subcorpus en
español, también de 50~documentos de temática equivalente procedentes de
revistas y obras de referencia hispanohablantes. Esta separación responde a
dos motivaciones complementarias: por un lado, la especificidad de los
recursos lingüísticos requeridos por cada lengua ---modelos de
\textsc{pos-tagging}, \textit{stoplists} y patrones morfosintácticos
diferenciados---; por otro, el interés contrastivo de obtener un inventario
que refleje las particularidades del léxico especializado en cada tradición
académica.

El proceso se articula en cuatro fases secuenciales (véase la Figura~\ref{fig:diagrama_extraccion}):
(i)~preprocesamiento y limpieza del corpus; (ii)~filtrado lingüístico mediante
patrones morfosintácticos; (iii)~ponderación estadística de los candidatos;
y (iv)~validación experta y organización de cada inventario. En las secciones a continuación se describe cada una de ellas.

\begin{figure}[ht] % 'ht' ayuda a que se ubique 'aquí' o en la parte 'superior'
    \centering
    \begin{tikzpicture}[
        node distance=1.2cm,
        auto,
        block/.style={rectangle, draw, fill=blue!5, text width=9cm, text centered, rounded corners, minimum height=0.8cm, font=\small},
        subblock/.style={rectangle, draw, fill=gray!5, text width=4cm, text centered, rounded corners, minimum height=0.8cm, font=\footnotesize},
        filter/.style={rectangle, draw, dashed, fill=orange!5, text width=8.5cm, font=\scriptsize},
        arrow/.style={-Stealth, thick}
    ]

    % --- ENTRADA ---
    \node (corpus_en) [subblock] {\textbf{Subcorpus Inglés} \\ 50 documentos};
    \node (corpus_es) [subblock, right=0.8cm of corpus_en] {\textbf{Subcorpus Español} \\ 50 documentos};

    % --- FASE 1 ---
    \node (fase1) [block, below=1cm of $(corpus_en.south)!0.5!(corpus_es.south)$] {
        \textbf{I. Preprocesamiento y Anotación} \\
        \scriptsize Limpieza (\textsc{pdf}), \texttt{langdetect} y modelos \texttt{spaCy} (\textit{lg/news})
    };

    % --- FASE 2 ---
    \node (fase2) [block, below=0.8cm of fase1] {
        \textbf{II. Filtrado Lingüístico} \\
        \scriptsize Aplicación de Patrones POS (\textit{NP, de/of, Nominales})
    };

    % --- FILTROS (Algoritmo 1) ---
    \node (filtros) [filter, below=0.4cm of fase2] {
        \textbf{Filtros de Calidad:} \\
        Longitud ($1<n<6$) $\cdot$ Categorial (funcionales) $\cdot$ Normalización $\cdot$ \textit{Stoplists} $\cdot$ Freq $\ge 2$
    };

    % --- FASE 3 ---
    \node (fase3) [block, below=0.4cm of filtros] {
        \textbf{III. Ponderación Estadística} \\
        \scriptsize Índice: $0,6 \cdot \widetilde{\text{TF-IDF}} + 0,4 \cdot \widetilde{\text{C-value}}$ (+ bono \textit{corpus})
    };

    % --- FASE 4 ---
    \node (fase4) [block, below=0.8cm of fase3] {
        \textbf{IV. Validación Experta} \\
        \scriptsize Cribado por percentil 90 y categorización temática
    };

    % --- SALIDA ---
    \node (inv_en) [subblock, below=1cm of fase4, xshift=-2.3cm, fill=green!5] {
        \textbf{Inventario EN} \\ 170 términos validados
    };
    \node (inv_es) [subblock, below=1cm of fase4, xshift=2.3cm, fill=green!5] {
        \textbf{Inventario ES} \\ 148 términos validados
    };

    % --- LÍNEAS ---
    \draw [arrow] (corpus_en.south) -- ++(0,-0.4) -| (fase1.north);
    \draw [arrow] (corpus_es.south) -- ++(0,-0.4) -| (fase1.north);
    \draw [arrow] (fase1) -- (fase2);
    \draw [arrow] (fase2) -- (filtros);
    \draw [arrow] (filtros) -- (fase3);
    \draw [arrow] (fase3) -- (fase4);
    \draw [arrow] (fase4.south) -- ++(0,-0.4) -| (inv_en.north);
    \draw [arrow] (fase4.south) -- ++(0,-0.4) -| (inv_es.north);

    \end{tikzpicture}
    
    \caption{Flujo secuencial de las fases de extracción terminológica.}
    \label{fig:diagrama_extraccion}
\end{figure}

Antes de la extracción propiamente dicha, los textos se sometieron a una
limpieza estructural orientada a eliminar el ruido tipográfico y documental
habitual en textos digitalizados o exportados desde \textsc{pdf}. Las
operaciones aplicadas incluyeron la supresión de las secciones de bibliografía
y referencias, la eliminación de pies de figura y tabla, la limpieza de
encabezados de sección aislados, la normalización de guiones de final de línea
(palabras partidas) y el descarte de líneas compuestas mayoritariamente por
cifras o marcadores numéricos de cita.

La detección automática de la lengua se efectuó sobre los primeros
3\,000~caracteres del texto limpio mediante la librería
\texttt{langdetect}\footnote{\url{https://pypi.org/project/langdetect/}} y
determinó el modelo de \texttt{spaCy} cargado en cada caso:
\texttt{en\_core\_web\_lg} para los documentos en inglés y
\texttt{es\_core\_news\_md} para los documentos en español. La anotación
morfosintáctica resultante (tokens, categorías gramaticales \textsc{pos} y
lemas) constituye la entrada de la fase siguiente. La Figura~\ref{fig:viz_spacy} ilustra el
resultado de esta anotación sobre dos fragmentos representativos del corpus: uno del
subcorpus en inglés y otro del subcorpus en español. Cada token se muestra con su
etiqueta \textsc{pos} (estándar Universal Dependencies); los fondos sombreados y los
marcadores identifican los tokens que forman unidades terminológicas validadas
(\texttt{B-T}: inicio de término; \texttt{I-T}: continuación).

\begin{figure}[H]
    \centering
    \includegraphics[width=\linewidth]{images/viz_spacy_pos.png}
    \caption{Anotación morfosintáctica de dos fragmentos del corpus TECLA mediante
             \texttt{spaCy} (Universal Dependencies). Panel superior: subcorpus EN
             (\textit{In corpus linguistics, parallel corpora are usually smaller\ldots});
             panel inferior: subcorpus ES
             (\textit{En la lematización y el etiquetado morfosintáctico de las
             palabras del corpus\ldots}). Los fondos en verde claro señalan tokens
             que forman parte de una unidad terminológica validada (\texttt{B-T}/\texttt{I-T}).
             La variedad de etiquetas \textsc{pos} en una misma oración muestra la
             riqueza del contexto sintáctico explotado por los patrones de extracción.}
    \label{fig:viz_spacy}
\end{figure}

La identificación de candidatos a término se basa en un conjunto de patrones
morfosintácticos formalizados sobre etiquetas \textsc{pos} y definidos
diferenciadamente por lengua. Los patrones recogen las estructuras
sintagmáticas más productivas en la formación del léxico especializado:
sintagmas nominales pre- y posnominales, secuencias de sustantivos y
construcciones mediadas por preposición (\textit{of}~/~\textit{de}). La
Tabla~\ref{tab:patrones-pos} sintetiza la gramática terminológica empleada
junto con ejemplos representativos.

\begin{table}[H]
\centering
\small
\caption{Gramática terminológica: patrones \textsc{pos} de extracción
         por lengua, con ejemplos representativos.}
\label{tab:patrones-pos}
\begin{tabular}{@{} c p{4.2cm} p{4.2cm} p{4.5cm} @{}}
\toprule
 & \textbf{ID} & \textbf{Patrón POS} & \textbf{Ejemplos} \\
\midrule
\multicolumn{4}{@{}l}{\textbf{Inglés (\texttt{en})}} \\
\midrule
1 & \texttt{EN\_NP\_PRENOMINAL}  & \texttt{(ADJ|NOUN|PROPN)* NOUN} & \textit{balanced corpus}, \textit{annotation scheme} \\
2 & \texttt{EN\_NP\_OF}          & \texttt{NOUN+ of NOUN+}          & \textit{part of speech}, \textit{analysis of variance} \\
3 & \texttt{EN\_NOUN\_SEQUENCE}  & \texttt{NOUN NOUN+}              & \textit{token frequency}, \textit{corpus benchmark} \\
4 & \texttt{EN\_NP\_PROPN}       & \texttt{PROPN PROPN+}            & \textit{British National Corpus} \\
\midrule
\multicolumn{4}{@{}l}{\textbf{Español (\texttt{es})}} \\
\midrule
5  & \texttt{ES\_NP\_PRENOMINAL}  & \texttt{(ADJ|NOUN|PROPN)* NOUN}  & \textit{patrón colocacional} \\
6  & \texttt{ES\_NP\_POSTNOMINAL} & \texttt{NOUN ADJ+}               & \textit{corpus balanceado}, \textit{frecuencia léxica} \\
7  & \texttt{ES\_NP\_DE}          & \texttt{NOUN+ de NOUN+}          & \textit{esquema de anotación}, \textit{lista de palabras} \\
8  & \texttt{ES\_NP\_DE\_ADJ}     & \texttt{NOUN+ de NOUN ADJ+}      & \textit{análisis de varianza estadística} \\
9  & \texttt{ES\_NP\_PROPN}       & \texttt{PROPN PROPN+}            & \textit{CREA Académico} \\
10 & \texttt{ES\_NP\_NOUN\_PROPN} & \texttt{NOUN PROPN+}             & \textit{corpus PRESEEA} \\
\bottomrule
\end{tabular}
\end{table}

Mientras que la Tabla~\ref{tab:patrones-pos} describe la gramática
formalizada \emph{ex ante}, la Tabla~\ref{tab:patrones_pos} muestra el
comportamiento empírico de dichos patrones tras la validación experta:
recoge la distribución observada de las estructuras morfosintácticas en el
inventario final del corpus.

\begin{table}[H]
\centering
\small
\setlength{\tabcolsep}{4pt}
\renewcommand{\arraystretch}{1.15}
\caption{Distribución empírica de patrones morfosintácticos en los términos
  validados del corpus \textsc{tecla} (EN: $n=155$ tipos; ES: $n=146$ tipos).
  Porcentajes calculados sobre el total de tipos únicos por lengua.}
\label{tab:patrones_pos}
\begin{tabular}{@{} l r r @{\quad} l r r @{}}
\toprule
\multicolumn{3}{c}{\textbf{Inglés (EN)}} & \multicolumn{3}{c}{\textbf{Español (ES)}} \\
\cmidrule(lr){1-3}\cmidrule(lr){4-6}
\textbf{Patrón} & \textbf{$n$} & \textbf{\%} &
\textbf{Patrón} & \textbf{$n$} & \textbf{\%} \\
\midrule
Uni-token (N/PROPN)      & 63 & 40{,}6 & Uni-token (N/PROPN)   & 95 & 65{,}1 \\
(ADJ|N|PROPN)* N         & 30 & 19{,}4 & N ADJ$^{+}$           & 18 & 12{,}3 \\
N N$^{+}$                & 27 & 17{,}4 & PROPN PROPN$^{+}$     & 16 & 11{,}0 \\
PROPN PROPN$^{+}$        & 16 & 10{,}3 & N$^{+}$ de N$^{+}$    &  6 &  4{,}1 \\
Compuesto con guion      &  8 &  5{,}2 & (ADJ|N|PROPN)* N      &  5 &  3{,}4 \\
Otros                    &  7 &  4{,}5 & Otros                 &  3 &  2{,}1 \\
N$^{+}$ of N$^{+}$       &  4 &  2{,}6 & N PROPN$^{+}$         &  2 &  1{,}4 \\
                         &    &        & N$^{+}$ de N ADJ$^{+}$&  1 &  0{,}7 \\
\midrule
\textbf{Total}           & \textbf{155} & \textbf{100} &
\textbf{Total}           & \textbf{146} & \textbf{100} \\
\bottomrule
\end{tabular}
\end{table}

Los datos confirman varias regularidades del léxico especializado del dominio.
En primer lugar, los uni-tokens dominan en ambas lenguas, con una asimetría
notable a favor del español (65,1\,\% frente al 40,6\,\% del inglés): la
mayor productividad morfológica del español permite lexicalizar mediante
derivación nominal (\textit{corpus}, \textit{anotación}, \textit{concordancia})
conceptos que en inglés requieren composición sintagmática. En segundo
lugar, las construcciones multipalabra revelan el contraste tipológico
esperado entre ambas lenguas: el inglés se apoya preferentemente en el
patrón prenominal \texttt{(ADJ|N|PROPN)* N} (19,4\,\%; \textit{natural
language processing}, \textit{inter-annotator agreement}), mientras que el
español privilegia el patrón posnominal \texttt{N ADJ}$^{+}$ (12,3\,\%;
\textit{lingüística computacional}, \textit{frecuencia léxica}). Por último,
las secuencias de nombres propios mantienen un peso comparable en ambas
lenguas (10-11\,\%) y corresponden, casi sin excepción, a denominaciones de
recursos y herramientas (\textit{British National Corpus},
\textit{WordSmith Tools}).

El procedimiento de extracción sigue el Algoritmo~\ref{alg:extraccion}: para
cada oración del corpus anotado se comprueba si alguna secuencia de tokens
satisface alguno de los patrones definidos para la lengua del documento;
en caso afirmativo, el candidato se somete a una batería de filtros de
calidad y, si los supera, se normaliza canónicamente y se incorpora al
inventario.

\begin{algorithm}[H]
\caption{Extracción de candidatos a término.}
\label{alg:extraccion}
\KwEntrada{Corpus $\mathcal{C} = \{d_1, \dots, d_n\}$, patrones POS $\mathcal{G}$,
           \textit{stoplists} por lengua $\mathcal{S}_\ell$,
           longitud máxima \texttt{max\_tokens}\,=\,6.}
\KwSalida{Lista de candidatos normalizados $\mathcal{T}$.}
\ParaCada{documento $d$ de $\mathcal{C}$}{
    Detectar la lengua del documento $\rightarrow \ell$\;
    Cargar el modelo \texttt{spaCy} para $\ell$ y anotar $d$
    (tokens, etiquetas \textsc{pos}, lemas) $\rightarrow \mathcal{A}$\;
    \ParaCada{oración $s$ de $\mathcal{A}$}{
        \ParaCada{patrón $p$ de $\mathcal{G}[\ell]$}{
            \ParaCada{secuencia $\sigma$ que coincide con $p$ en $s$}{
                \lSi{$|\sigma|=0$ \textbf{o} $|\sigma|>\texttt{max\_tokens}$}{\textbf{descartar}}
                \lSi{$\sigma$ contiene determinantes, pronombres u otras palabras funcionales}{\textbf{descartar}}
                Construir la forma canónica $\lambda$ concatenando los lemas en minúscula\;
                \lSi{forma superficial de $\sigma$ o $\lambda \in \mathcal{S}_\ell$}{\textbf{descartar}}
                \lSi{$|\sigma|=1$ \textbf{y} $\lambda$ es \textit{stopword}}{\textbf{descartar}}
                Añadir $(\sigma,\lambda,d)$ a $\mathcal{T}$\;
            }
        }
    }
}
Consolidar $\mathcal{T}$: acumular frecuencias absolutas y documentales\;
\lSi{frecuencia$(\lambda) < 2$ \textbf{o} doc\_freq$(\lambda) < 1$}{\textbf{descartar}}
\Devolver{$\mathcal{T}$ ordenada por forma canónica.}
\end{algorithm}

La cadena de filtros responde a cinco criterios complementarios. El
\textit{filtro de longitud} descarta las secuencias vacías y las que
superan los seis tokens. El \textit{filtro categorial} elimina los
candidatos que contienen determinantes, pronombres u otras palabras
funcionales. La \textit{normalización canónica} reduce las secuencias
supervivientes a su forma lematizada en minúsculas. El \textit{filtro de
\textit{stoplist}} rechaza el candidato cuando su forma superficial o
canónica coincide con alguna entrada de la lista de exclusión académica
por lengua ---construida a partir de los recursos de \texttt{spaCy} y
\texttt{NLTK} y complementada con los ficheros específicos
\texttt{stop\_en}/\texttt{stop\_es}---; en el caso de los unigramas, el
rechazo se aplica también cuando el lema coincide con una \textit{stopword}.
Por último, el \textit{filtro de frecuencia mínima} ---aplicado en la fase
de consolidación--- descarta los candidatos con frecuencia absoluta inferior
a~2 o presentes en un único documento, por considerarse estadísticamente
insuficientes para su validación.

El análisis estructural de los candidatos supervivientes revela una
regularidad marcada: en ambos subcorpus, la mayoría corresponde a
secuencias formadas por sustantivos consecutivos, seguidas a distancia
por las construcciones preposicionales y los sintagmas adjetivo-sustantivo.
Este predominio de la composición nominal confirma empíricamente que la
terminología de la lingüística de corpus se lexicaliza preferentemente
mediante compuestos nominales, en consonancia con lo descrito en la
bibliografía sobre terminología científica y técnica
\autocite{cabre1999terminology, sager1990terminology}

Una vez delimitado el conjunto de candidatos por criterios lingüísticos, se
aplicaron dos métricas complementarias para estimar el grado de
\emph{terminologicidad} de cada unidad: \textit{TF-IDF} y \textit{C-value}
(véase~§\ref{subsec:ate_metodos}). Ambas se reescalaron al intervalo $[0,1]$
mediante normalización mínimo-máximo y se combinaron en un índice
compuesto mediante ponderación lineal:

\begin{equation}
    \text{terminologicidad}(t) \;=\; 0{,}6 \cdot \widetilde{\text{TF-IDF}}(t)
                                \;+\; 0{,}4 \cdot \widetilde{\text{C-value}}(t).
    \label{eq:terminologicidad}
\end{equation}

A este índice se incorpora una bonificación cualitativa para los candidatos
que contienen explícitamente el núcleo léxico \textit{corpus} ---o su
equivalente en español---, en reconocimiento de su papel estructurante en
el dominio y sin alterar la ponderación relativa de las métricas. La
justificación empírica de esta decisión puede leerse a la luz del análisis de
incrustación terminológica presentado en §\ref{subsec:terminos_incrustados}: el
término \textit{corpus} actúa como núcleo de 28~cadenas terminológicas en el
subcorpus inglés y de 26~en el español, con tasas de incrustación del 32\,\%
y el 13\,\% respectivamente; su frecuencia bruta sobrestima sistemáticamente su
terminologicidad autónoma, y la bonificación compensa parcialmente el descuento
que el C-value aplica a sus ocurrencias incrustadas. El
índice resultante permite ordenar los candidatos a lo largo de un
\textit{continuum} de especialización y aplicar umbrales de corte
---por ejemplo, el percentil~90--- para concentrar la validación experta
en aquellos con mayor probabilidad de constituir términos válidos. Los
resultados muestran un comportamiento diferenciado y complementario de
ambas métricas: el \textit{C-value} alcanza valores más altos en las
primeras posiciones del \textit{ranking} para unidades poliléxicas con
fuerte cohesión sintagmática, mientras que el \textit{TF-IDF} discrimina
mejor a lo largo del \textit{continuum} al penalizar las unidades con
distribución demasiado difusa.

Los candidatos situados en el percentil superior del índice de
terminologicidad fueron sometidos a validación experta con un doble
objetivo: confirmar su pertinencia como unidades terminológicas del
dominio y asignarles una categoría temática. El criterio de validación
combinó los principios formales descritos en la metodología (véase
Capítulo~\ref{chap:3_extraccion}) con el conocimiento experto de la
disciplina, prestando especial atención a la estabilidad denominativa,
la autonomía referencial y la pertenencia al léxico nuclear de la
lingüística de corpus.

El proceso arrojó dos inventarios validados: 170~términos en el subcorpus
inglés, distribuidos en 11~categorías temáticas, y 148~términos en el
subcorpus español, agrupados en 6~macrocategorías. El mayor peso
proporcional de los unigramas en el subcorpus español (cf.
Tabla~\ref{tab:patrones_pos}) puede explicarse, como ya se anticipó, por
la productividad morfológica de la lengua, que permite lexicalizar
conceptos especializados mediante derivación nominal sin necesidad de
recurrir sistemáticamente a compuestos sintagmáticos. Las
Figuras~\ref{fig:categorias_en} y~\ref{fig:categorias_es} ilustran la
distribución por categorías en cada subcorpus.

\begin{figure}[H]
    \centering
    \includegraphics[width=0.8\linewidth]{images/fig_categorias_en.png}
    \caption{Distribución de los términos validados del subcorpus en
             inglés por categoría temática.}
    \label{fig:categorias_en}
\end{figure}

\begin{figure}[H]
    \centering
    \includegraphics[width=0.8\linewidth]{images/fig_categorias_en.png}
    \caption{Distribución de los términos validados del subcorpus en
             español por categoría temática (equivalente a EN).}
    \label{fig:categorias_es}
\end{figure}

En el subcorpus inglés predominan los términos vinculados al
\textit{diseño y compilación de corpus} (34~unidades), seguidos por
\textit{frecuencia y estadística} (25) y \textit{léxico y fraseología}
(25). En el subcorpus español, la categoría \textit{recursos y corpus}
(\textsc{rec-cor}) concentra el mayor número de unidades (52), seguida
por \textit{herramientas y software} (\textsc{her-soft}, 41~términos) y
\textit{metodología y diseño} (\textsc{met-di}, 26~términos). Esta última
es directamente comparable con la categoría \textit{diseño y compilación
de corpus} del subcorpus inglés, lo que indica una cobertura temática
equivalente para el núcleo metodológico del dominio. La diferencia más
notable reside en el peso de \textsc{rec-cor} en español, que excede
considerablemente la proporción de su equivalente inglés ---hecho
coherente con la mayor atención que la bibliografía hispanohablante
dedica a la descripción y catalogación de corpus disponibles y
herramientas de procesamiento, frecuentemente articulada como revisión
historiográfica de proyectos nacionales (\textit{CREA}, \textit{CORDE},
\textit{CORPES~XXI}, \textit{PRESEEA}).

\begin{figure}[H]
    \centering
    \includegraphics[width=\linewidth]{images/fig_categorias_comparativa.png}
    \caption{Comparación normalizada de la distribución por macrocategorías
             entre el subcorpus en inglés y el subcorpus en español.}
    \label{fig:categorias_comparativa}
\end{figure}

\subsubsection{Términos vs. Entidades: Una distinción tipológica}
\label{subsec:terminos_vs_entidades}

El análisis del inventario de unidades terminológicas validadas pone de manifiesto una heterogeneidad tipológica que trasciende la distinción convencional entre términos y no-términos. Aunque el conjunto total de 318~unidades (170~en inglés, 148~en español) comparte un valor especializado indiscutible y potencial lexicográfico como material de base para la construcción de glosarios bilingües, no todas ellas constituyen términos en sentido estricto.

La distinción que proponemos en las tablas siguientes separa dos categorías de unidades complementarias:

\begin{enumerate}
    \item \textbf{Términos:} unidades léxicas que designan conceptos, procesos, propiedades o metodologías especializadas del dominio (ej.\ \textit{representatividad}, \textit{anotación semántica}, \textit{frecuencia relativa}). Son las unidades que responden a la caracterización formal, estadística y distribucional desarrollada en la sección anterior.

    \item \textbf{Entidades:} denominaciones propias de recursos, proyectos, herramientas y referencias concretas (ej.\ \textit{British National Corpus}, \textit{Sketch Engine}, \textit{CREA}). Poseen un marcado valor léxico-referencial para la construcción de glosarios, pero designan particulares institucionales o comerciales antes que conceptos generalizables.
\end{enumerate}

Esta separación tipológica revela un contraste lingüístico y académico profundo entre las dos tradiciones bibliográficas (Tabla~\ref{tab:terminos_vs_entidades}): el subcorpus inglés concentra su terminología en conceptos especializados (84{,}7\,\% términos, 15{,}3\,\% entidades), reflejando una tradición de investigación teórica consolidada y orientada hacia la formalización de conceptos; el subcorpus español, en cambio, dedica una proporción sustancialmente mayor de su vocabulario a la documentación de recursos, proyectos y herramientas de origen local (37{,}2\,\% términos, 62{,}8\,\% entidades), evidenciando un énfasis historiográfico y un mayor detalle en la descripción y catalogación de infraestructuras nacionales.

\begin{table}[H]
\centering
\small
\caption{Distribución de unidades terminológicas por tipo lingüístico.
         \textit{Unidades}: total de tipos validados.
         \textit{Términos}: unidades con valor conceptual especializado (stricto sensu).
         \textit{Entidades}: denominaciones de recursos, herramientas y proyectos con valor lexicográfico.}
\label{tab:terminos_vs_entidades}
\begin{tabular}{@{} l r r r r r r @{}}
\toprule
\textbf{Lengua} & \textbf{Unidades totales} & \multicolumn{2}{c}{\textbf{Términos}} & \multicolumn{2}{c}{\textbf{Entidades}} \\
\cmidrule(lr){3-4}\cmidrule(lr){5-6}
 & & $n$ & \% & $n$ & \% \\
\midrule
Inglés (EN) & 170 & 144 & 84{,}7\,\% & 26 & 15{,}3\,\% \\
Español (ES) & 148 & 55 & 37{,}2\,\% & 93 & 62{,}8\,\% \\
\midrule
\textbf{Total} & \textbf{318} & \textbf{199} & \textbf{62{,}6\,\%} & \textbf{119} & \textbf{37{,}4\,\%} \\
\bottomrule
\end{tabular}
\end{table}

Las Tablas~\ref{tab:terminos_en} y~\ref{tab:terminos_es} recogen el
inventario completo de términos validados para cada subcorpus junto con
sus valores estadísticos asociados. El inventario inglés
(Tabla~\ref{tab:terminos_en}) evidencia la madurez de la disciplina en
esta lengua a través de una terminología altamente estratificada: las
unidades con mayor índice de terminologicidad no son únicamente las más
frecuentes (\textit{corpus}, \textit{annotation}), sino aquellas que
operan como nodos conceptuales de los que derivan subdominios enteros.
El inventario captura simultáneamente la terminología técnica del
procesamiento (\textit{POS tagging}, \textit{lemmatisation}) y la de
base estadística (\textit{log-likelihood}, \textit{confidence interval}),
lo que demuestra que el corpus de trabajo cubre satisfactoriamente las
vertientes teórica y aplicada de la lingüística de corpus anglosajona.

\small
\begin{longtable}{@{} r p{5cm} r r r r @{}}
\toprule
\# & \textbf{UT} & \textbf{Frec.\ abs.} & \textbf{DF} & \textbf{TF--IDF} & \textbf{C-value} \\
\midrule
\endfirsthead
\multicolumn{6}{l}{\small\itshape (continuación)} \\
\toprule
\# & \textbf{UT} & \textbf{Frec.\ abs.} & \textbf{DF} & \textbf{TF--IDF} & \textbf{C-value} \\
\midrule
\endhead
\midrule
\multicolumn{6}{r}{\small\itshape Continúa en la página siguiente\ldots} \\
\endfoot
\bottomrule
\caption{Términos validados del subcorpus en inglés, organizados por
         categoría temática.}
\label{tab:terminos_en}
\endlastfoot
\multicolumn{6}{l}{\textbf{Conceptos fundacionales}} \\
\midrule
1. & \textit{corpus} & 13767 & 50 & 1{,}0000 & 1{,}0000 \\
2. & \textit{type} & 1048 & 41 & 0{,}0974 & 0{,}1888 \\
3. & \textit{corpus linguistics} & 2138 & 47 & 0{,}1689 & 0{,}1558 \\
4. & \textit{corpus-based} & 416 & 38 & 0{,}0419 & 0{,}0397 \\
5. & \textit{corpus-driven} & 31 & 9 & 0{,}0075 & 0{,}0030 \\
\midrule
\multicolumn{6}{l}{\textbf{Diseño y compilación de corpus}} \\
\midrule
6. & \textit{representativeness} & 555 & 22 & 0{,}0866 & 0{,}0415 \\
7. & \textit{speech corpus} & 151 & 6 & 0{,}0424 & 0{,}0310 \\
8. & \textit{text type} & 266 & 24 & 0{,}0392 & 0{,}0486 \\
9. & \textit{corpus design} & 319 & 20 & 0{,}0528 & 0{,}0426 \\
10. & \textit{genre} & 297 & 33 & 0{,}0342 & 0{,}0396 \\
11. & \textit{sampling frame} & 91 & 5 & 0{,}0270 & 0{,}0158 \\
12. & \textit{corpus size} & 104 & 13 & 0{,}0216 & 0{,}0155 \\
13. & \textit{reference corpus} & 113 & 20 & 0{,}0187 & 0{,}0190 \\
14. & \textit{national corpus} & 189 & 26 & 0{,}0263 & 0{,}0185 \\
15. & \textit{general corpus} & 85 & 7 & 0{,}0227 & 0{,}0106 \\
16. & \textit{corpus building} & 93 & 7 & 0{,}0248 & 0{,}0095 \\
17. & \textit{multimodal corpus} & 33 & 3 & 0{,}0112 & 0{,}0062 \\
18. & \textit{learner corpus} & 64 & 9 & 0{,}0156 & 0{,}0081 \\
19. & \textit{monitor corpus} & 42 & 7 & 0{,}0112 & 0{,}0067 \\
20. & \textit{training corpus} & 31 & 6 & 0{,}0087 & 0{,}0056 \\
21. & \textit{subcorpus} & 84 & 16 & 0{,}0158 & 0{,}0048 \\
22. & \textit{web corpus} & 32 & 6 & 0{,}0090 & 0{,}0045 \\
23. & \textit{comparable corpus} & 38 & 7 & 0{,}0101 & 0{,}0063 \\
24. & \textit{annotated corpus} & 70 & 10 & 0{,}0163 & 0{,}0060 \\
25. & \textit{multilingual corpus} & 41 & 8 & 0{,}0104 & 0{,}0050 \\
26. & \textit{corpus compilation} & 43 & 15 & 0{,}0083 & 0{,}0058 \\
27. & \textit{spoken corpus} & 63 & 15 & 0{,}0122 & 0{,}0055 \\
28. & \textit{monolingual corpus} & 30 & 3 & 0{,}0102 & 0{,}0031 \\
29. & \textit{diachronic corpus} & 29 & 11 & 0{,}0065 & 0{,}0047 \\
30. & \textit{mega corpora} & 39 & 3 & 0{,}0132 & 0{,}0025 \\
31. & \textit{historical corpus} & 19 & 10 & 0{,}0044 & 0{,}0042 \\
32. & \textit{specialized corpus} & 27 & 12 & 0{,}0058 & 0{,}0033 \\
33. & \textit{balanced corpus} & 31 & 10 & 0{,}0072 & 0{,}0030 \\
34. & \textit{web crawling} & 5 & 3 & 0{,}0017 & 0{,}0006 \\
35. & \textit{synchronic} & 94 & 16 & 0{,}0176 & 0{,}0090 \\
36. & \textit{parallel corpus} & 145 & 8 & 0{,}0369 & 0{,}0138 \\
37. & \textit{written corpus} & 15 & 9 & 0{,}0037 & 0{,}0014 \\
38. & \textit{sample size} & 217 & 9 & 0{,}0528 & 0{,}0207 \\
39. & \textit{stratified sampling} & 5 & 1 & 0{,}0021 & 0{,}0001 \\
\midrule
\multicolumn{6}{l}{\textbf{Frecuencia y estadística}} \\
\midrule
40. & \textit{frequency} & 1667 & 40 & 0{,}1592 & 0{,}1822 \\
41. & \textit{sample} & 1127 & 31 & 0{,}1369 & 0{,}1553 \\
42. & \textit{probability} & 380 & 17 & 0{,}0691 & 0{,}0438 \\
43. & \textit{range} & 584 & 37 & 0{,}0604 & 0{,}0539 \\
44. & \textit{dispersion} & 329 & 16 & 0{,}0618 & 0{,}0312 \\
45. & \textit{occurrence} & 308 & 24 & 0{,}0454 & 0{,}0423 \\
46. & \textit{population} & 286 & 16 & 0{,}0537 & 0{,}0287 \\
47. & \textit{association measures} & 102 & 6 & 0{,}0286 & 0{,}0146 \\
48. & \textit{descriptive statistics} & 55 & 8 & 0{,}0140 & 0{,}0047 \\
49. & \textit{relative frequency} & 28 & 9 & 0{,}0068 & 0{,}0052 \\
50. & \textit{raw frequency} & 10 & 3 & 0{,}0034 & 0{,}0033 \\
51. & \textit{mutual information} & 31 & 9 & 0{,}0075 & 0{,}0031 \\
52. & \textit{inferential statistics} & 22 & 6 & 0{,}0062 & 0{,}0017 \\
53. & \textit{absolute frequency} & 5 & 2 & 0{,}0018 & 0{,}0014 \\
54. & \textit{chi-squared} & 46 & 8 & 0{,}0117 & 0{,}0003 \\
55. & \textit{log-likelihood} & 39 & 12 & 0{,}0084 & 0{,}0037 \\
56. & \textit{standard deviation} & 121 & 14 & 0{,}0243 & 0{,}0115 \\
57. & \textit{statistical significance} & 34 & 10 & 0{,}0079 & 0{,}0032 \\
58. & \textit{confidence interval} & 46 & 6 & 0{,}0129 & 0{,}0044 \\
59. & \textit{t-score} & 15 & 3 & 0{,}0051 & 0{,}0014 \\
60. & \textit{mean score} & 77 & 2 & 0{,}0285 & 0{,}0073 \\
61. & \textit{z-score} & 1 & 1 & 0{,}0004 & 0{,}0001 \\
62. & \textit{deviation of proportions} & 2 & 2 & 0{,}0007 & 0{,}0003 \\
63. & \textit{normalized frequency} & 10 & 5 & 0{,}0030 & 0{,}0000 \\
64. & \textit{effect size} & 60 & 8 & 0{,}0153 & 0{,}0002 \\
\midrule
\multicolumn{6}{l}{\textbf{Anotación y marcado}} \\
\midrule
65. & \textit{annotation} & 1513 & 38 & 0{,}1524 & 0{,}1450 \\
66. & \textit{annotation scheme} & 126 & 16 & 0{,}0237 & 0{,}0230 \\
67. & \textit{tagset} & 155 & 13 & 0{,}0323 & 0{,}0161 \\
68. & \textit{word class} & 37 & 8 & 0{,}0094 & 0{,}0168 \\
69. & \textit{accuracy} & 165 & 29 & 0{,}0212 & 0{,}0150 \\
70. & \textit{corpus annotation} & 132 & 18 & 0{,}0232 & 0{,}0136 \\
71. & \textit{recall} & 107 & 17 & 0{,}0194 & 0{,}0053 \\
72. & \textit{POS tagging} & 72 & 11 & 0{,}0161 & 0{,}0085 \\
73. & \textit{treebank} & 105 & 13 & 0{,}0219 & 0{,}0039 \\
74. & \textit{parts of speech} & 75 & 16 & 0{,}0141 & 0{,}0099 \\
75. & \textit{annotation tools} & 51 & 8 & 0{,}0130 & 0{,}0049 \\
76. & \textit{markup language} & 40 & 9 & 0{,}0097 & 0{,}0044 \\
77. & \textit{syntactic annotation} & 32 & 7 & 0{,}0085 & 0{,}0028 \\
78. & \textit{semantic annotation} & 34 & 12 & 0{,}0073 & 0{,}0029 \\
79. & \textit{gold standard} & 13 & 5 & 0{,}0039 & 0{,}0012 \\
80. & \textit{morphological annotation} & 7 & 2 & 0{,}0026 & 0{,}0006 \\
81. & \textit{part-of-speech tagging} & 77 & 15 & 0{,}0149 & 0{,}0073 \\
82. & \textit{inter-annotator agreement} & 21 & 5 & 0{,}0062 & 0{,}0020 \\
83. & \textit{inter-rater reliability} & 1 & 1 & 0{,}0004 & 0{,}0001 \\
84. & \textit{inline annotation} & 2 & 1 & 0{,}0008 & 0{,}0002 \\
\midrule
\multicolumn{6}{l}{\textbf{Léxico y fraseología}} \\
\midrule
85. & \textit{keyword} & 337 & 16 & 0{,}0633 & 0{,}0622 \\
86. & \textit{collocation} & 260 & 19 & 0{,}0444 & 0{,}0318 \\
87. & \textit{collocate} & 95 & 9 & 0{,}0231 & 0{,}0237 \\
88. & \textit{frequency list} & 201 & 15 & 0{,}0390 & 0{,}0355 \\
89. & \textit{grammar} & 423 & 29 & 0{,}0542 & 0{,}0331 \\
90. & \textit{keyword analysis} & 139 & 12 & 0{,}0300 & 0{,}0262 \\
91. & \textit{word list} & 96 & 14 & 0{,}0193 & 0{,}0156 \\
92. & \textit{word frequency} & 61 & 15 & 0{,}0118 & 0{,}0155 \\
93. & \textit{function word} & 88 & 13 & 0{,}0183 & 0{,}0142 \\
94. & \textit{hapax legomena} & 67 & 7 & 0{,}0179 & 0{,}0062 \\
95. & \textit{content words} & 38 & 10 & 0{,}0089 & 0{,}0065 \\
96. & \textit{keyness} & 86 & 8 & 0{,}0219 & 0{,}0036 \\
97. & \textit{lexical bundles} & 47 & 5 & 0{,}0140 & 0{,}0043 \\
98. & \textit{node word} & 18 & 5 & 0{,}0053 & 0{,}0033 \\
99. & \textit{semantic prosody} & 29 & 3 & 0{,}0098 & 0{,}0027 \\
100. & \textit{colligation} & 15 & 4 & 0{,}0047 & 0{,}0019 \\
101. & \textit{collocational strength} & 6 & 4 & 0{,}0019 & 0{,}0006 \\
102. & \textit{stopword} & 4 & 2 & 0{,}0015 & 0{,}0004 \\
103. & \textit{phraseology} & 13 & 7 & 0{,}0035 & 0{,}0006 \\
104. & \textit{formulaic language} & 8 & 4 & 0{,}0025 & 0{,}0005 \\
105. & \textit{n-gram} & 37 & 6 & 0{,}0104 & 0{,}0035 \\
106. & \textit{type-token ratio} & 11 & 3 & 0{,}0037 & 0{,}0010 \\
107. & \textit{lexico-grammar} & 4 & 3 & 0{,}0014 & 0{,}0004 \\
108. & \textit{semantic preference} & 2 & 1 & 0{,}0008 & 0{,}0003 \\
109. & \textit{multiword expression} & 1 & 1 & 0{,}0004 & 0{,}0003 \\
\midrule
\multicolumn{6}{l}{\textbf{Procesamiento del lenguaje}} \\
\midrule
110. & \textit{token} & 279 & 20 & 0{,}0462 & 0{,}0385 \\
111. & \textit{alignment} & 153 & 16 & 0{,}0287 & 0{,}0139 \\
112. & \textit{tagger} & 142 & 22 & 0{,}0222 & 0{,}0153 \\
113. & \textit{dataset} & 122 & 25 & 0{,}0175 & 0{,}0159 \\
114. & \textit{regular expression} & 148 & 11 & 0{,}0332 & 0{,}0135 \\
115. & \textit{stem} & 40 & 10 & 0{,}0093 & 0{,}0072 \\
116. & \textit{lemma} & 89 & 10 & 0{,}0208 & 0{,}0082 \\
117. & \textit{lexeme} & 40 & 5 & 0{,}0119 & 0{,}0065 \\
118. & \textit{clustering} & 73 & 4 & 0{,}0231 & 0{,}0049 \\
119. & \textit{parsing} & 54 & 17 & 0{,}0098 & 0{,}0038 \\
120. & \textit{lemmatisation} & 34 & 5 & 0{,}0101 & 0{,}0029 \\
121. & \textit{tokenisation} & 25 & 5 & 0{,}0074 & 0{,}0020 \\
122. & \textit{word sense disambiguation} & 15 & 4 & 0{,}0047 & 0{,}0039 \\
123. & \textit{segmentation} & 28 & 14 & 0{,}0056 & 0{,}0020 \\
124. & \textit{text mining} & 12 & 6 & 0{,}0034 & 0{,}0017 \\
125. & \textit{dependency parsing} & 5 & 4 & 0{,}0016 & 0{,}0009 \\
126. & \textit{named entity recognition} & 3 & 3 & 0{,}0010 & 0{,}0005 \\
\midrule
\multicolumn{6}{l}{\textbf{Corpus de referencia}} \\
\midrule
127. & \textit{BNC} & 815 & 22 & 0{,}1272 & 0{,}0222 \\
128. & \textit{Brown Corpus} & 215 & 16 & 0{,}0404 & 0{,}0210 \\
129. & \textit{LOB Corpus} & 141 & 8 & 0{,}0359 & 0{,}0113 \\
130. & \textit{British National Corpus} & 114 & 22 & 0{,}0178 & 0{,}0171 \\
131. & \textit{Helsinki Corpus} & 45 & 9 & 0{,}0110 & 0{,}0045 \\
132. & \textit{Penn Treebank} & 41 & 9 & 0{,}0100 & 0{,}0037 \\
133. & \textit{GloWbE} & 37 & 4 & 0{,}0117 & 0{,}0015 \\
134. & \textit{LOB} & 304 & 11 & 0{,}0682 & 0{,}0002 \\
135. & \textit{International Corpus of English} & 51 & 14 & 0{,}0102 & 0{,}0097 \\
136. & \textit{COHA} & 67 & 9 & 0{,}0163 & 0{,}0064 \\
137. & \textit{Corpus of Contemporary American English} & 46 & 9 & 0{,}0112 & 0{,}0102 \\
138. & \textit{Bank of English} & 47 & 5 & 0{,}0140 & 0{,}0071 \\
139. & \textit{MICASE} & 30 & 5 & 0{,}0089 & 0{,}0029 \\
140. & \textit{BNCWeb} & 188 & 4 & 0{,}0594 & 0{,}0179 \\
141. & \textit{COCA} & 291 & 13 & 0{,}0606 & 0{,}0001 \\
\midrule
\multicolumn{6}{l}{\textbf{Metadatos y formatos}} \\
\midrule
142. & \textit{XML} & 567 & 23 & 0{,}0860 & 0{,}0302 \\
143. & \textit{metadata} & 251 & 24 & 0{,}0370 & 0{,}0170 \\
144. & \textit{header} & 94 & 7 & 0{,}0251 & 0{,}0097 \\
145. & \textit{encoding} & 143 & 16 & 0{,}0269 & 0{,}0109 \\
146. & \textit{HTML} & 221 & 27 & 0{,}0299 & 0{,}0027 \\
147. & \textit{plain text} & 42 & 10 & 0{,}0098 & 0{,}0038 \\
148. & \textit{UTF-8} & 75 & 6 & 0{,}0211 & 0{,}0015 \\
149. & \textit{SGML} & 64 & 9 & 0{,}0156 & 0{,}0010 \\
150. & \textit{TEI} & 60 & 12 & 0{,}0130 & 0{,}0057 \\
151. & \textit{unicode} & 59 & 10 & 0{,}0138 & 0{,}0001 \\
\midrule
\multicolumn{6}{l}{\textbf{Búsqueda y consulta}} \\
\midrule
152. & \textit{query} & 422 & 19 & 0{,}0720 & 0{,}0460 \\
153. & \textit{concordance} & 161 & 18 & 0{,}0284 & 0{,}0198 \\
154. & \textit{concordance lines} & 16 & 6 & 0{,}0045 & 0{,}0028 \\
155. & \textit{wildcard} & 16 & 2 & 0{,}0059 & 0{,}0018 \\
156. & \textit{KWIC} & 41 & 6 & 0{,}0115 & 0{,}0039 \\
157. & \textit{corpus queries} & 45 & 5 & 0{,}0134 & 0{,}0043 \\
\midrule
\multicolumn{6}{l}{\textbf{Herramientas y software}} \\
\midrule
158. & \textit{CLAWS} & 106 & 13 & 0{,}0221 & 0{,}0074 \\
159. & \textit{concordancer} & 21 & 6 & 0{,}0059 & 0{,}0033 \\
160. & \textit{Sketch Engine} & 25 & 8 & 0{,}0064 & 0{,}0021 \\
161. & \textit{WordSmith Tools} & 9 & 6 & 0{,}0025 & 0{,}0006 \\
162. & \textit{TreeTagger} & 59 & 7 & 0{,}0157 & 0{,}0056 \\
163. & \textit{CQPweb} & 21 & 7 & 0{,}0056 & 0{,}0020 \\
164. & \textit{Perl} & 70 & 6 & 0{,}0197 & 0{,}0067 \\
165. & \textit{CWB} & 6 & 5 & 0{,}0018 & 0{,}0006 \\
166. & \textit{NLTK} & 7 & 3 & 0{,}0024 & 0{,}0007 \\
167. & \textit{Wmatrix} & 6 & 3 & 0{,}0020 & 0{,}0006 \\
168. & \textit{AntConc} & 138 & 9 & 0{,}0336 & 0{,}0000 \\
\midrule
\multicolumn{6}{l}{\textbf{Disciplinas afines}} \\
\midrule
169. & \textit{natural language processing} & 90 & 24 & 0{,}0133 & 0{,}0137 \\
170. & \textit{computational linguistics} & 136 & 26 & 0{,}0189 & 0{,}0095 \\
\end{longtable}
\normalsize

El inventario español (Tabla~\ref{tab:terminos_es}) ofrece, en cambio, un
perfil parcialmente distinto. El núcleo metodológico es convergente con
el del inglés ---se replican unidades como \textit{representatividad},
\textit{corpus de referencia} o \textit{corpus paralelo}---, pero la
categoría \textsc{rec-cor} muestra una dispersión mucho mayor, alimentada
por la atención sostenida que la bibliografía hispanohablante presta a
la revisión historiográfica y descriptiva de proyectos nacionales
(\textit{CREA}, \textit{CORDE}, \textit{CORPES~XXI}, \textit{PRESEEA},
\textit{IULA}). Asimismo, los valores de \textit{C-value} revelan una
particularidad estructural: la dependencia de estructuras preposicionales
complejas (\texttt{N + de + N}) para la formación de términos
equivalentes a los compuestos nominales compactos del inglés se traduce
en una longitud media de candidato ligeramente superior. Conviene notar,
por último, que la escala absoluta de \textit{TF-IDF} y \textit{C-value}
en esta tabla difiere de la del subcorpus inglés porque los valores no
fueron renormalizados al intervalo $[0,1]$ tras la consolidación del
inventario español; las comparaciones interlingüísticas, por tanto, deben
realizarse en términos relativos (posición en el \textit{ranking})
y no absolutos.

\small
\begin{longtable}{@{} r p{5cm} r r r r @{}}
\toprule
\# & \textbf{UT} & \textbf{Frec.\ abs.} & \textbf{DF} & \textbf{TF--IDF} & \textbf{C-value} \\
\midrule
\endfirsthead
\multicolumn{6}{l}{\small\itshape (continuación)} \\
\toprule
\# & \textbf{UT} & \textbf{Frec.\ abs.} & \textbf{DF} & \textbf{TF--IDF} & \textbf{C-value} \\
\midrule
\endhead
\midrule
\multicolumn{6}{r}{\small\itshape Continúa en la página siguiente\ldots} \\
\endfoot
\bottomrule
\caption{Términos validados del subcorpus en español, organizados por
         categoría temática.}
\label{tab:terminos_es}
\endlastfoot
\multicolumn{6}{l}{\textbf{Metodología y diseño (\textsc{met-di})}} \\
\midrule
1. & \textit{corpus} & 11749 & 50 & 11749{,}0000 & 11691{,}5769 \\
2. & \textit{lingüística de corpus} & 690 & 47 & 750{,}3494 & 1073{,}0196 \\
3. & \textit{corpus de referencia} & 195 & 31 & 326{,}1229 & 309{,}0677 \\
4. & \textit{representatividad} & 223 & 34 & 344{,}1207 & 223{,}0000 \\
5. & \textit{corpus paralelo} & 130 & 17 & 325{,}3250 & 130{,}0000 \\
6. & \textit{corpus diacrónico} & 96 & 26 & 184{,}0836 & 96{,}0000 \\
7. & \textit{equilibrio} & 96 & 23 & 200{,}3964 & 96{,}0000 \\
8. & \textit{corpus de aprendices} & 41 & 12 & 121{,}8514 & 64{,}9835 \\
9. & \textit{web as corpus} & 40 & 8 & 140{,}1000 & 63{,}3985 \\
10. & \textit{corpus comparables} & 40 & 12 & 118{,}8794 & 40{,}0000 \\
11. & \textit{monitor} & 44 & 13 & 126{,}0631 & 33{,}0000 \\
12. & \textit{corpus especializado} & 32 & 12 & 95{,}1035 & 32{,}0000 \\
13. & \textit{corpus multilingüe} & 27 & 11 & 83{,}3615 & 27{,}0000 \\
14. & \textit{lingüística de corpus computacional} & 13 & 2 & 66{,}1370 & 26{,}0000 \\
15. & \textit{corpus general} & 25 & 11 & 77{,}1866 & 25{,}0000 \\
16. & \textit{corpus virtual} & 23 & 9 & 77{,}0614 & 23{,}0000 \\
17. & \textit{corpus bilingüe} & 16 & 7 & 58{,}7588 & 16{,}0000 \\
18. & \textit{unidad de muestra} & 9 & 1 & 51{,}0518 & 14{,}2647 \\
19. & \textit{corpus monitor} & 11 & 6 & 42{,}5158 & 11{,}0000 \\
20. & \textit{ETBC} & 10 & 2 & 50{,}8746 & 10{,}0000 \\
21. & \textit{corpus sincrónico} & 10 & 7 & 36{,}7243 & 10{,}0000 \\
22. & \textit{megacorpus} & 8 & 3 & 37{,}3794 & 8{,}0000 \\
23. & \textit{web for corpus} & 5 & 2 & 25{,}4373 & 7{,}9248 \\
24. & \textit{corpus telemático} & 5 & 1 & 28{,}3621 & 5{,}0000 \\
25. & \textit{corpus oportunista} & 3 & 3 & 14{,}0173 & 3{,}0000 \\
26. & \textit{corpus estático} & 2 & 2 & 10{,}1749 & 2{,}0000 \\
\midrule
\multicolumn{6}{l}{\textbf{Recursos y corpus (\textsc{rec-cor})}} \\
\midrule
27. & \textit{CREA} & 377 & 36 & 551{,}5404 & 377{,}0000 \\
28. & \textit{CORDE} & 324 & 27 & 604{,}2828 & 324{,}0000 \\
29. & \textit{CDH} & 104 & 12 & 309{,}0865 & 104{,}0000 \\
30. & \textit{PRESEEA} & 99 & 14 & 273{,}7879 & 99{,}0000 \\
31. & \textit{CAES} & 90 & 8 & 315{,}2250 & 90{,}0000 \\
32. & \textit{CORPES XXI} & 76 & 15 & 203{,}1043 & 76{,}0000 \\
33. & \textit{IULA} & 69 & 12 & 205{,}0670 & 69{,}0000 \\
34. & \textit{Brown Corpus} & 54 & 15 & 144{,}3110 & 54{,}0000 \\
35. & \textit{Val.Es.Co} & 51 & 11 & 157{,}4606 & 51{,}0000 \\
36. & \textit{CEMC} & 44 & 10 & 141{,}3717 & 44{,}0000 \\
37. & \textit{BNC} & 40 & 12 & 118{,}8794 & 40{,}0000 \\
38. & \textit{WordNet} & 39 & 5 & 159{,}4111 & 39{,}0000 \\
39. & \textit{COBUILD} & 35 & 15 & 93{,}5349 & 35{,}0000 \\
40. & \textit{CHARTA} & 33 & 10 & 106{,}0288 & 33{,}0000 \\
41. & \textit{CORLEC} & 32 & 10 & 102{,}8158 & 32{,}0000 \\
42. & \textit{Survey of English Usage} & 16 & 9 & 53{,}6080 & 32{,}0000 \\
43. & \textit{CEDEL2} & 27 & 9 & 90{,}4634 & 27{,}0000 \\
44. & \textit{CODEA+ 2015} & 24 & 4 & 104{,}4119 & 24{,}0000 \\
45. & \textit{C-ORAL-ROM} & 21 & 8 & 73{,}5525 & 21{,}0000 \\
46. & \textit{OPUS} & 19 & 6 & 73{,}4363 & 19{,}0000 \\
47. & \textit{CLEAE} & 15 & 1 & 85{,}0864 & 15{,}0000 \\
48. & \textit{Europarl} & 15 & 6 & 57{,}9761 & 15{,}0000 \\
49. & \textit{esTenTen18} & 14 & 5 & 57{,}2245 & 14{,}0000 \\
50. & \textit{SPLLOC} & 14 & 5 & 57{,}2245 & 14{,}0000 \\
51. & \textit{MUST} & 13 & 7 & 47{,}7415 & 13{,}0000 \\
52. & \textit{ODE} & 10 & 2 & 50{,}8746 & 10{,}0000 \\
53. & \textit{London-Lund Corpus} & 10 & 6 & 38{,}6507 & 10{,}0000 \\
54. & \textit{Survey of Spoken English} & 5 & 4 & 21{,}7525 & 10{,}0000 \\
55. & \textit{ACTRES} & 9 & 3 & 42{,}0518 & 9{,}0000 \\
56. & \textit{SPT} & 9 & 2 & 45{,}7872 & 9{,}0000 \\
57. & \textit{CORLEXIN} & 9 & 2 & 45{,}7872 & 9{,}0000 \\
58. & \textit{MULTINOT} & 7 & 3 & 32{,}7070 & 7{,}0000 \\
59. & \textit{ECPC} & 6 & 2 & 30{,}5248 & 6{,}0000 \\
60. & \textit{GENTT} & 6 & 4 & 26{,}1030 & 6{,}0000 \\
61. & \textit{MELLANGE} & 6 & 2 & 30{,}5248 & 6{,}0000 \\
62. & \textit{BIDTEA} & 6 & 2 & 30{,}5248 & 6{,}0000 \\
63. & \textit{COHA} & 5 & 3 & 23{,}3621 & 5{,}0000 \\
64. & \textit{CATE} & 5 & 4 & 21{,}7525 & 5{,}0000 \\
65. & \textit{LETRAC} & 4 & 2 & 20{,}3499 & 4{,}0000 \\
66. & \textit{CdE-Web} & 3 & 1 & 17{,}0173 & 3{,}0000 \\
67. & \textit{CORANE} & 3 & 2 & 15{,}2624 & 3{,}0000 \\
68. & \textit{Global Voices} & 3 & 2 & 15{,}2624 & 3{,}0000 \\
69. & \textit{COCA} & 3 & 3 & 14{,}0173 & 3{,}0000 \\
70. & \textit{CdE-NOW} & 2 & 1 & 11{,}3449 & 2{,}0000 \\
71. & \textit{COWS-L2H} & 2 & 1 & 11{,}3449 & 2{,}0000 \\
72. & \textit{FonoELE} & 2 & 2 & 10{,}1749 & 2{,}0000 \\
73. & \textit{GloWbE} & 2 & 2 & 10{,}1749 & 2{,}0000 \\
74. & \textit{LANGSNAP} & 2 & 1 & 11{,}3449 & 2{,}0000 \\
75. & \textit{CORESPI} & 1 & 1 & 5{,}6724 & 1{,}0000 \\
76. & \textit{ELEACTAR} & 1 & 1 & 5{,}6724 & 1{,}0000 \\
77. & \textit{Lancaster-Oslo-Bergen Corpus} & 1 & 1 & 5{,}6724 & 1{,}0000 \\
78. & \textit{Macrocorpus sociolingüístico} & 1 & 1 & 5{,}6724 & 1{,}0000 \\
\midrule
\multicolumn{6}{l}{\textbf{Herramientas y software (\textsc{her-soft})}} \\
\midrule
79. & \textit{FreeLing} & 119 & 13 & 340{,}9434 & 119{,}0000 \\
80. & \textit{Sketch Engine} & 113 & 16 & 292{,}1008 & 113{,}0000 \\
81. & \textit{BootCaT} & 61 & 5 & 249{,}3352 & 61{,}0000 \\
82. & \textit{TreeTagger} & 59 & 9 & 197{,}6793 & 59{,}0000 \\
83. & \textit{AntConc} & 36 & 10 & 115{,}6678 & 36{,}0000 \\
84. & \textit{Gensim} & 32 & 2 & 162{,}7988 & 32{,}0000 \\
85. & \textit{WordSmith Tools} & 28 & 6 & 108{,}2220 & 28{,}0000 \\
86. & \textit{GRAMPAL} & 24 & 2 & 122{,}0991 & 24{,}0000 \\
87. & \textit{IMS Open Corpus Workbench} & 12 & 2 & 61{,}0496 & 24{,}0000 \\
88. & \textit{Stanza} & 20 & 2 & 101{,}7493 & 20{,}0000 \\
89. & \textit{SDL Trados Studio} & 11 & 1 & 62{,}3967 & 17{,}4346 \\
90. & \textit{MALLET} & 16 & 2 & 81{,}3994 & 16{,}0000 \\
91. & \textit{Trados} & 24 & 5 & 98{,}0991 & 13{,}0000 \\
92. & \textit{TagAnt} & 11 & 3 & 51{,}3967 & 11{,}0000 \\
93. & \textit{taggers} & 11 & 3 & 51{,}3967 & 11{,}0000 \\
94. & \textit{CQP} & 10 & 2 & 50{,}8746 & 10{,}0000 \\
95. & \textit{IXA} & 8 & 2 & 40{,}6997 & 8{,}0000 \\
96. & \textit{Xbench} & 8 & 1 & 45{,}3794 & 8{,}0000 \\
97. & \textit{rtweet} & 6 & 1 & 34{,}0346 & 6{,}0000 \\
98. & \textit{Wmatrix} & 6 & 3 & 28{,}0346 & 6{,}0000 \\
99. & \textit{LF Aligner} & 5 & 3 & 23{,}3621 & 5{,}0000 \\
100. & \textit{T-Hoarder} & 5 & 1 & 28{,}3621 & 5{,}0000 \\
101. & \textit{NLTK} & 4 & 4 & 17{,}4020 & 4{,}0000 \\
102. & \textit{Transit} & 5 & 3 & 23{,}3621 & 3{,}0000 \\
103. & \textit{Bitext2tmx} & 3 & 2 & 15{,}2624 & 3{,}0000 \\
104. & \textit{ChoRus} & 3 & 1 & 17{,}0173 & 3{,}0000 \\
105. & \textit{Gephi} & 3 & 1 & 17{,}0173 & 3{,}0000 \\
106. & \textit{ConcGram} & 2 & 1 & 11{,}3449 & 2{,}0000 \\
107. & \textit{DMI-TCAT} & 2 & 1 & 11{,}3449 & 2{,}0000 \\
108. & \textit{GDEX} & 2 & 1 & 11{,}3449 & 2{,}0000 \\
109. & \textit{GenEx} & 2 & 1 & 11{,}3449 & 2{,}0000 \\
110. & \textit{MultiTrans} & 2 & 2 & 10{,}1749 & 2{,}0000 \\
111. & \textit{Star Transit} & 2 & 2 & 10{,}1749 & 2{,}0000 \\
112. & \textit{Tweepy} & 2 & 1 & 11{,}3449 & 2{,}0000 \\
113. & \textit{Wordfast} & 2 & 2 & 10{,}1749 & 2{,}0000 \\
114. & \textit{align} & 2 & 2 & 10{,}1749 & 2{,}0000 \\
115. & \textit{AntPConc} & 2 & 2 & 10{,}1749 & 2{,}0000 \\
116. & \textit{Notepad++} & 1 & 1 & 5{,}6724 & 1{,}0000 \\
117. & \textit{Pencil} & 1 & 1 & 5{,}6724 & 1{,}0000 \\
118. & \textit{Stanford Parser} & 1 & 1 & 5{,}6724 & 1{,}0000 \\
119. & \textit{Stingray} & 1 & 1 & 5{,}6724 & 1{,}0000 \\
\midrule
\multicolumn{6}{l}{\textbf{Procesamiento lingüístico (\textsc{proc-lin})}} \\
\midrule
120. & \textit{lematización} & 101 & 23 & 210{,}8337 & 101{,}0000 \\
121. & \textit{alineación} & 97 & 14 & 268{,}2569 & 97{,}0000 \\
122. & \textit{concordancia} & 69 & 23 & 144{,}0349 & 69{,}0000 \\
123. & \textit{anotación semántica} & 35 & 8 & 122{,}5875 & 35{,}0000 \\
124. & \textit{transcriptor} & 26 & 2 & 132{,}2740 & 26{,}0000 \\
125. & \textit{etiquetado morfosintáctico} & 25 & 8 & 87{,}5625 & 25{,}0000 \\
126. & \textit{lingüística computacional de corpus} & 12 & 2 & 61{,}0496 & 24{,}0000 \\
127. & \textit{tokenización} & 13 & 6 & 50{,}2459 & 13{,}0000 \\
128. & \textit{anotación pragmática} & 10 & 5 & 40{,}8746 & 10{,}0000 \\
129. & \textit{KWIC} & 9 & 8 & 31{,}5225 & 9{,}0000 \\
130. & \textit{UKB} & 6 & 2 & 30{,}5248 & 6{,}0000 \\
131. & \textit{LESK} & 4 & 2 & 20{,}3499 & 4{,}0000 \\
132. & \textit{desambiguación morfosintáctica} & 3 & 3 & 14{,}0173 & 3{,}0000 \\
\midrule
\multicolumn{6}{l}{\textbf{Estadística y léxico (\textsc{est-lex})}} \\
\midrule
133. & \textit{keywords} & 45 & 23 & 93{,}9358 & 45{,}0000 \\
134. & \textit{frecuencia absoluta} & 35 & 10 & 112{,}4548 & 35{,}0000 \\
135. & \textit{n-gramas} & 33 & 8 & 115{,}5825 & 33{,}0000 \\
136. & \textit{TTR} & 31 & 3 & 144{,}8452 & 31{,}0000 \\
137. & \textit{frecuencia relativa} & 27 & 8 & 94{,}5675 & 27{,}0000 \\
138. & \textit{ley de Zipf} & 11 & 4 & 47{,}8555 & 17{,}4346 \\
139. & \textit{medias móviles} & 14 & 1 & 79{,}4140 & 14{,}0000 \\
140. & \textit{word sketch} & 13 & 5 & 53{,}1370 & 13{,}0000 \\
141. & \textit{wordlist} & 7 & 3 & 32{,}7070 & 7{,}0000 \\
142. & \textit{índice normalizado de dispersión} & 3 & 1 & 17{,}0173 & 6{,}0000 \\
143. & \textit{lista de frecuencia} & 3 & 3 & 14{,}0173 & 4{,}7549 \\
144. & \textit{frecuencia por millón de palabras} & 2 & 1 & 11{,}3449 & 4{,}6439 \\
145. & \textit{keyness} & 2 & 2 & 10{,}1749 & 2{,}0000 \\
\midrule
\multicolumn{6}{l}{\textbf{Tecnología y formatos (\textsc{tec-for})}} \\
\midrule
146. & \textit{EAGLES} & 38 & 11 & 117{,}3236 & 38{,}0000 \\
147. & \textit{Universal Dependencies} & 6 & 3 & 28{,}0346 & 6{,}0000 \\
148. & \textit{Universal POS} & 2 & 2 & 10{,}1749 & 2{,}0000 \\
\end{longtable}
\normalsize

% ============================================================
% SECCIÓN: Modelado supervisado
% ============================================================

\section{Modelado supervisado de unidades terminológicas}
\label{sec:modelo_tc}

La identificación automática de unidades terminológicas puede abordarse desde dos perspectivas complementarias: la primera, descrita en la sección precedente (§\ref{sec:extraccion}), opera a nivel de \emph{tipo} léxico, ponderando estadísticamente los candidatos extraídos mediante reglas lingüísticas; la segunda, que se desarrolla en el presente apartado, opera a nivel de \emph{token}, tratando cada ocurrencia individual como una instancia de clasificación en un marco supervisado. Mientras que el enfoque híbrido permite construir un inventario de candidatos validados, el modelo supervisado responde a una pregunta diferente y metodológicamente más exigente: ¿en qué medida pueden aprenderse automáticamente, a partir de rasgos lingüísticos y estadísticos, los patrones que distinguen un término de una unidad léxica no especializada?

\subsection{Clasificadores token-based: limitaciones para ATE}
\label{subsec:clasificadores_tokenbase}

El enfoque más intuitivo para la detección de términos es tratar cada token como una instancia de clasificación binaria (\textit{término} vs.\ \textit{no-término}) mediante algoritmos estándar de aprendizaje supervisado: regresión logística, bosques aleatorios (Random Forest), máquinas de vectores de soporte (SVM) o redes neuronales densas. Estos clasificadores operan sobre vectores de rasgos (features) asociados a cada token ---categoría gramatical, frecuencia, contexto distribucional, etc.--- y predicen una etiqueta de clase.

\paragraph{Por qué parecen funcionar.}
En un esquema de validación cruzada (cross-validation) tradicional sobre datos token-anotados, estos clasificadores token-based pueden alcanzar valores de F1 de 0{,}70--0{,}80 o incluso superiores. El motivo es que muchos tokens denominan la \emph{misma} unidad terminológica a lo largo del corpus: si la palabra ``corpus'' aparece 13\,767 veces en el subcorpus inglés, un clasificador que haya visto cientos de ocurrencias de ``corpus'' etiquetadas como término en el conjunto de entrenamiento puede memorizar simplemente el patrón \textit{palabra = ``corpus''} $\rightarrow$ \textit{etiqueta = término}, sin necesidad de aprender rasgos contextuales más generales.

Este fenómeno de \emph{memorización léxica} inflama artificialmente las métricas de evaluación. La validación cruzada tradicional divide el corpus en folds \emph{aleatoriamente a nivel de token}, lo que significa que la \textbf{misma palabra} (y frecuentemente la misma unidad terminológica multi-token) aparece distribuida entre folds de entrenamiento y de test. Un clasificador que memoriza la identidad léxica obtiene F1 alto porque los mismos tipos terminológicos se repiten entre train y test.

\paragraph{Por qué fallan para terminología.}
El problema es que la métrica de evaluación no responde la pregunta correcta. Un sistema que ``aprende'' que ``corpus'' es un término porque memoriza esa palabra específica fracasa cuando se enfrenta a un término que nunca ha visto: la tarea real de extracción terminológica automática es detectar términos \emph{nuevos} o poco frecuentes, no reetiquetizar las palabras más frecuentes del dominio.

Este es el punto crítico: la evaluación de clasificadores token-based sobre datos sin control de leakage no distingue entre (i) aprendizaje genuino de patrones contextuales que generalizan a términos nuevos, y (ii) memorización léxica de identidades de palabras. Para tareas de \emph{automatic term extraction}, donde el objetivo es detectar términos desconocidos con precisión, la memorización es una falla silenciosa que métricas estándar no capturan.

\subsection{Clasificadores secuenciales: Conditional Random Fields (CRF)}
\label{subsec:crf}

Para capturar la naturaleza \emph{secuencial} de las unidades terminológicas (especialmente las multi-token), empleamos Conditional Random Fields (CRF), un modelo discriminativo para etiquetado de secuencias que modela las transiciones entre etiquetas de forma explícita. Un CRF asigna a cada token una etiqueta del esquema BIO (\texttt{B-T}: comienzo de término, \texttt{I-T}: interior de término, \texttt{O}: fuera de término), teniendo en cuenta no solo los rasgos del token individual sino también las restricciones de secuencia válidas (ej.\ una etiqueta \texttt{I-T} debe ir precedida de \texttt{B-T} o de otra \texttt{I-T}).

\paragraph{Extracción de rasgos para CRF.}
Para cada token se extraen aproximadamente 25 rasgos lingüísticos, estadísticos y contextuales:

\begin{enumerate}
    \item \textbf{Rasgos léxicos:} palabra en minúscula, lema, tres primeros caracteres, tres últimos caracteres.
    \item \textbf{Rasgos morfosintácticos:} etiqueta POS, es capitalizado, es todo mayúsculas, es todo minúsculas, es dígito, longitud.
    \item \textbf{Rasgos contextuales:} tokens anteriores/siguientes (ventana de contexto $\pm 2$), etiquetas POS adyacentes.
    \item \textbf{Rasgos estadísticos:} frecuencia bruta en el corpus, frecuencia de documento, diversidad contextual (número de contextos discursivos distintos).
\end{enumerate}

El vector de rasgos se construye dinámicamente para cada ocurrencia del token, permitiendo que el modelo capture tanto propiedades invariantes de la palabra (ej.\ su categoría gramatical típica) como propiedades variables del contexto (ej.\ qué tipo de palabra lo precede).

\paragraph{Evaluación convencional: validación cruzada sin control de leakage.}
En una primera serie de experimentos, entrenamos un CRF sobre datos divididos en 5 folds mediante validación cruzada estratificada. Los resultados fueron aparentemente satisfactorios:

\begin{itemize}
    \item \textbf{EN:} F1 token-level promedio $\approx$ 0{,}67 $\pm$ 0{,}04
    \item \textbf{ES:} F1 token-level promedio $\approx$ 0{,}73 $\pm$ 0{,}03
\end{itemize}

Sin embargo, estos números ocultan un sesgo metodológico grave: aunque los folds se dividían aleatoriamente, \textbf{los términos no se dividían}. Una unidad terminológica como ``corpus linguistics'' podía aparecer en múltiples ocurrencias distribuidas entre train y test. El modelo podía, en principio, memorizar que la secuencia \texttt{(NOUN, NOUN)} en contextos específicos es un término, sin aprender patrones distribucionales que generalizaran a términos totalmente nuevos.

\paragraph{Evaluación rigurosa: holdout a nivel de términos (span-level).}
Para diagnosticar si el CRF genuinamente aprendía o memorizaba, diseñamos un protocolo de evaluación \emph{span-level}:

\begin{enumerate}
    \item Extraer todas las instancias de términos únicos del corpus (ej.\ ``corpus linguistics'' como unidad léxica).
    \item Estratificar los términos por frecuencia: raros (<10 ocurrencias), medios (10--50), comunes (>50).
    \item Seleccionar un porcentaje de términos para holdout (10\%, 20\%, 30\%, 50\%).
    \item \textbf{Excluir del conjunto de entrenamiento TODAS las oraciones que contienen esos términos holdout}.
    \item Evaluar en las oraciones excluidas.
\end{enumerate}

Este protocolo elimina completamente el riesgo de que un término aparezca en ambos train y test, forzando al modelo a detectar términos basándose únicamente en \emph{rasgos distribucionales e patrones de contexto}, no en memorización léxica.

\paragraph{Resultados con holdout estricto.}
Los resultados fueron dramáticos y revelan el alcance del problema de memorización:

\begin{table}[H]
\centering
\small
\caption{Rendimiento de CRF con holdout span-level. F1 token-level sobre conjuntos de test generados por exclusión de oraciones con términos holdout. Resultados incluyen intervalos de confianza 95\%.}
\label{tab:crf_holdout_results}
\begin{tabular}{@{} l r r r r @{}}
\toprule
\textbf{Lengua} & \textbf{Holdout \%} & \textbf{Tokens Test} & \textbf{F1 Token-level} & \textbf{IC 95\%} \\
\midrule
\multirow{4}{*}{\textbf{EN}} & 10\% & 284{,}748 & 0{,}1564 & $\pm 0{,}0013$ \\
 & 20\% & 296{,}407 & 0{,}3406 & $\pm 0{,}0015$ \\
 & 30\% & 299{,}187 & 0{,}3338 & $\pm 0{,}0011$ \\
 & 50\% & 301{,}019 & 0{,}6761 & $\pm 0{,}0017$ \\
\midrule
\multirow{4}{*}{\textbf{ES}} & 10\% & 359{,}369 & 0{,}4775 & $\pm 0{,}0016$ \\
 & 20\% & 682{,}589 & 0{,}7569 & $\pm 0{,}0012$ \\
 & 30\% & 685{,}263 & 0{,}7881 & $\pm 0{,}0010$ \\
 & 50\% & 710{,}325 & 0{,}7694 & $\pm 0{,}0009$ \\
\bottomrule
\end{tabular}
\end{table}

El patrón es claro: cuando los términos holdout son pocos (10\%), el F1 se colapsa a $\approx 0{,}15$--0{,}48. A medida que el holdout aumenta (50\%), el F1 se recupera ($\approx 0{,}67$--0{,}77). ¿Por qué? Porque con holdout del 50\%, los términos restantes en test son suficientemente frecuentes como para que el modelo haya visto instancias similares en train. Con 10\% de holdout, los términos test son exactamente aquellos que el modelo nunca ha visto: el F1 bajo refleja un fracaso genuino de generalización.

\subsection{El problema de Zipfian tail: por qué CRF no puede aprender}
\label{subsec:zipfian_tail}

Los resultados de CRF plantean una pregunta fundamental: ¿es el problema inherente al algoritmo, o es un artefacto del dominio? Para responder, analizamos la distribución de frecuencias de términos en el corpus (Figura~\ref{fig:zipf_terms}).

\begin{figure}[H]
    \centering
    \includegraphics[width=0.8\linewidth]{experimento_8_julio/graficas/03_zipf_law_TERMINOS.png}
    \caption{Distribución de Zipf de TÉRMINOS a nivel de tipos únicos. Panel izquierdo: inglés (209 términos únicos, 5{,}355 instancias); panel derecho: español (115 términos únicos, 12{,}424 instancias). Escala log-log mostrando ajuste teórico de Zipf.}
    \label{fig:zipf_terms}
\end{figure}

\begin{figure}[H]
    \centering
    \includegraphics[width=0.8\linewidth]{experimento_8_julio/graficas/04_heaps_law_TERMINOS.png}
    \caption{Ley de Heaps para TÉRMINOS: crecimiento de vocabulario vs. tamaño del corpus. Panel izquierdo: inglés (β = 0{,}393); panel derecho: español (β = 0{,}591). Escala log-log mostrando saturación terminológica temprana.}
    \label{fig:heaps_terms}
\end{figure}

La Ley de Zipf describe una distribución que aparece sistemáticamente en lenguaje natural: la frecuencia de un tipo es inversamente proporcional a su rango. En nuestro corpus:

\begin{itemize}
    \item \textbf{EN:} El término más frecuente (``corpus'') aparece 1{,}366 veces; la mediana es 4 ocurrencias; el 24{,}4\,\% de términos únicos aparecen \textbf{solo 1 vez}.
    \item \textbf{ES:} El término más frecuente (``corpus'') aparece 7{,}949 veces; la mediana es 5 ocurrencias; el 19{,}1\,\% de términos únicos aparecen \textbf{solo 1 vez}.
\end{itemize}

Este es el corazón del problema: con una distribución Zipfiana, es \textbf{imposible} que un modelo estadístico estándar (como CRF) memorice suficientemente todos los términos. Un término que aparece una sola vez no ofrece al modelo entrenando suficientes instancias para aprender características discriminativas de su identidad léxica.

Más aún: incluso un término que aparece 5 veces (mediana) ofrece muy pocas instancias de aprendizaje. Un modelo basado en patrones léxicos debe fracasar cuando la mayoría de sus datos de prueba corresponde precisamente a ese tail Zipfiano ---términos raros con pocas o ningunas apariciones en entrenamiento.

La solución, entonces, no es mejor optimización del CRF, sino un cambio fundamental en lo que el modelo \emph{aprende}. En lugar de memorizar identidades léxicas, el modelo debe aprender \textbf{patrones distribucionales y contextuales} que generalicen a términos completamente nuevos. Esto es lo que hace un modelo basado en embeddings distribucionales: BiLSTM-CRF.

\subsection{BiLSTM-CRF: aprendizaje distribucional}
\label{subsec:bilstm_crf}

Un BiLSTM-CRF (Bidirectional Long Short-Term Memory con CRF en la capa de salida) reemplaza los rasgos manuales por representaciones continuas aprendidas:

\begin{enumerate}
    \item Cada token se codifica como un embedding (vector denso de dimensionalidad típicamente 100--300).
    \item Un LSTM bidireccional lee la secuencia hacia adelante y hacia atrás, produciendo una representación contextual de cada token que captura relaciones a largo plazo en la secuencia.
    \item Una capa CRF en la salida modela transiciones entre etiquetas (BIO).
\end{enumerate}

El punto crítico es que \textbf{los embeddings se aprenden durante el entrenamiento}. Dos palabras con contextos similares reciben embeddings similares, incluso si nunca aparecieron en los mismos contextos concretos. Un término raro puede beneficiarse de la información distribucional de palabras relacionadas: si el modelo aprendió que ``linguistic processing'' es un término, y ve a ``term extraction'' en un contexto distribucionalmente similar, puede inferir que también es un término, \emph{sin haber visto esa secuencia exacta en entrenamiento}.

\paragraph{Resultados de BiLSTM-CRF con holdout estricto.}
Al entrenar el mismo protocolo span-level holdout con BiLSTM-CRF:

\begin{table}[H]
\centering
\small
\caption{Rendimiento de BiLSTM-CRF con holdout span-level. F1 token-level sobre los mismos conjuntos de test que CRF (Tabla~\ref{tab:crf_holdout_results}).}
\label{tab:bilstm_holdout_results}
\begin{tabular}{@{} l r r r r @{}}
\toprule
\textbf{Lengua} & \textbf{Holdout \%} & \textbf{Tokens Test} & \textbf{F1 Token-level} & \textbf{IC 95\%} \\
\midrule
\multirow{4}{*}{\textbf{EN}} & 10\% & 284{,}748 & 0{,}8451 & $\pm 0{,}0008$ \\
 & 20\% & 296{,}407 & 0{,}8637 & $\pm 0{,}0007$ \\
 & 30\% & 299{,}187 & 0{,}8512 & $\pm 0{,}0009$ \\
 & 50\% & 301{,}019 & 0{,}8248 & $\pm 0{,}0011$ \\
\midrule
\multirow{4}{*}{\textbf{ES}} & 10\% & 359{,}369 & 0{,}9648 & $\pm 0{,}0005$ \\
 & 20\% & 682{,}589 & 0{,}9531 & $\pm 0{,}0006$ \\
 & 30\% & 685{,}263 & 0{,}9487 & $\pm 0{,}0007$ \\
 & 50\% & 710{,}325 & 0{,}9312 & $\pm 0{,}0009$ \\
\bottomrule
\end{tabular}
\end{table}

El contraste es dramático: BiLSTM-CRF mantiene F1 de 0{,}82--0{,}96 incluso con 10\% de holdout, donde CRF colapsa a 0{,}15--0{,}48. La diferencia es estructural: BiLSTM aprende patrones distribucionales que generalizan a términos nunca vistos; CRF, restringido a rasgos manuales, no puede escapar del tail Zipfiano.

\paragraph{Validación externa: ACTER cross-domain.}
Para confirmar que estos resultados no son artefactos del corpus específico, evaluamos el modelo BiLSTM-CRF entrenado en el corpus TECLA sobre el benchmark ACTER, que contiene anotaciones de términos en cuatro dominios distintos (wind energy, cardiovascular medicine, physics, quantitative semantics) en inglés.

Los resultados (Tabla~\ref{tab:acter_results}) muestran que el rendimiento es comparable al holdout span-level, validando que el aprendizaje es genuinamente transferible a dominios nuevos:

\begin{table}[H]
\centering
\small
\caption{Rendimiento de BiLSTM-CRF en ACTER (validación externa cross-domain). F1 token-level sobre cada dominio del benchmark ACTER.}
\label{tab:acter_results}
\begin{tabular}{@{} l r r @{}}
\toprule
\textbf{Dominio ACTER} & \textbf{F1 Promedio} & \textbf{IC 95\%} \\
\midrule
Wind Energy & 0{,}7812 & $\pm 0{,}0019$ \\
Cardiovascular Medicine & 0{,}7654 & $\pm 0{,}0021$ \\
Physics & 0{,}7423 & $\pm 0{,}0024$ \\
Quantitative Semantics & 0{,}7501 & $\pm 0{,}0022$ \\
\midrule
\textbf{Promedio ACTER} & \textbf{0{,}7598} & $\pm 0{,}0021$ \\
\bottomrule
\end{tabular}
\end{table}

El rendimiento promedio en ACTER (0{,}76) es ligeramente inferior al del holdout span-level en TECLA (0{,}82--0{,}96), pero la diferencia es coherente con la dificultad de transferencia cross-domain: dominios nuevos introducen vocabulario desconocido y patrones sintácticos distintos. Lo importante es que el modelo no colapsa: BiLSTM-CRF mantiene generalización robusta incluso en dominios completamente nuevos, confirmando que ha aprendido patrones lingüísticos genuinos, no memorización léxica del corpus de entrenamiento.

\paragraph{Conclusión.}
El análisis de Zipf, el protocolo span-level holdout y la validación external ACTER convergen en una conclusión clara: para tareas de extracción automática de terminología donde el objetivo es detectar términos nuevos y contextos nuevos, se requiere un modelo que aprenda patrones distribucionales, no identidades léxicas. BiLSTM-CRF proporciona este aprendizaje mediante embeddings contextuales, alcanzando generalización robusta incluso bajo evaluación rigurosa (span-level holdout) y validando en dominios completamente nuevos (ACTER).

\subsection{Importancia de rasgos: Rasgos contextuales vs. frecuenciales en la detección de términos no vistos}
\label{subsec:importancia_rasgos}

Más allá de las métricas agregadas de rendimiento, el análisis de la contribución relativa de cada rasgo al aprendizaje del modelo revela qué propiedades lingüísticas y estadísticas son más discriminativas para identificar un término, especialmente ante términos raros o completamente nuevos. Este análisis es crucial porque conecta directamente con un debate fundamental en la literatura de extracción automática de términos (ATE): ¿predicen mejor los rasgos contextuales (dispersión, patrones de co-ocurrencia, embeddings) o los rasgos de frecuencia (TF-IDF, C-value) la presencia de términos nuevos no vistos durante entrenamiento?

\paragraph{Antecedentes en la literatura: Rasgos contextuales y generalización a términos no vistos.}
La evidencia empírica en ATE sugiere, aunque de manera mayormente indirecta, que rasgos contextuales superan a rasgos puramente frecuenciales en la detección de términos raros y no vistos. Rigouts Terryn et al. (2021--2022), en su análisis comparativo sobre el corpus ACTER (48.4\,\% de términos únicos son hapax), evaluaron un CRF feature-based (177 features en 6 grupos) frente a una RNN con BiLSTM-CRF+BERT. En el escenario más realista (entrenar out-of-domain, testear in-domain), el recall en hapax fue 42{,}8\,\% para la RNN contextual frente a 34{,}9\,\% para el sistema de features (una reducción de la brecha hapax-overall a la mitad: −7{,}1 pp para RNN vs.\ −13{,}0 pp para HAMLET). En TermEval 2020 (shared task con dominio de test held-out, \textit{heart failure}), el sistema ganador TALN-LS2N, basado en BERT sin ninguna métrica estadística tradicional (TF-IDF, C-value), alcanzó 64\,\% de recall en hapax, mientras que un sistema de umbral de frecuencia obtuvo 0\,\%. La conclusión de los organizadores fue que TALN-LS2N "obtiene un recall sorprendentemente estable y alto (64\,\%) en términos hapax, probablemente porque no usa ninguna de las métricas estadísticas relacionadas con frecuencia".

Sin embargo, la evidencia cuantitativa en ATE es \textbf{mayormente indirecta}: típicamente se reporta recall estratificado por frecuencia (hapax vs.\ frecuentes) o evaluación cross-domain, no una partición explícita visto/no-visto con F1 separado. El test más directo de esta partición proviene de NER (Taillé et al., 2020), donde el gap seen-unseen con BiLSTM-CRF+BERT fue ≈13{,}5 puntos de F1, validando que la contextualización mitiga pero no elimina el problema de términos nuevos.

\paragraph{Novedad metodológica: El test LOTO explícito de esta tesis.}
Esta investigación introduce un protocolo menos común en ATE: partición explícita a nivel de TÉRMINOS (span-level holdout, ya descrito en §\ref{subsec:crf}) que genera conjuntos seen y unseen con F1 reportado por partición. Este diseño es más directo que los proxies habituales (hapax, cross-domain) y permite cuantificar cómo cada familia de rasgos predice términos completamente nuevos. A nuestro conocimiento, no existe precedente publicado de un test LOTO con este diseño en ATE supervisado.

\paragraph{Metodología de análisis de importancia de rasgos.}
Para los modelos CRF, la importancia de rasgos se cuantifica mediante los coeficientes de peso aprendidos durante el entrenamiento (rasgos con pesos absolutos mayores son más discriminativos). Para BiLSTM-CRF, se utiliza el método de permutación (permutation importance): se entrena el modelo completo, se perturba iterativamente cada rasgo, y se mide el descenso en F1; rasgos cuya perturbación causa mayor descenso son más importantes.

En esta investigación, el rasgo \textbf{diversidad contextual (CD)} se operacionaliza como el \textbf{número de documentos/oraciones discursivas distintos} en que aparece una unidad. Un documento es una unidad discursiva independiente (artículo, párrafo, sección); esta operacionalización captura la \emph{dispersión} de la unidad a través del dominio especializado. Cabe notar que CD es una \textbf{feature contextual artesanal} (handcrafted), conceptualmente diferente de los embeddings contextuales aprendidos (BERT, ELMo, BiLSTM) que predominan en la literatura reciente. La mayor parte de la evidencia cuantitativa a favor de rasgos contextuales procede de embeddings contextuales aprendidos, no de métricas simples como CD; por tanto, el valor de CD como predictor es un hallazgo parcial y específico de este corpus/dominio.

\begin{table}[H]
\centering
\small
\caption{Importancia de rasgos en BiLSTM-CRF para inglés. Ordenados por descenso de F1 tras perturbación. \textit{CD} = diversidad contextual; \textit{DF} = frecuencia de documento; \textit{TF} = frecuencia absoluta.}
\label{tab:bio_importancia_en}
\begin{tabular}{@{} l r r @{}}
\toprule
\textbf{Rasgo} & \textbf{Δ F1 (permutación)} & \textbf{\% Importancia} \\
\midrule
Diversidad contextual (CD) & 0{,}0847 & 28{,}5\,\% \\
Token anterior (POS) & 0{,}0623 & 20{,}9\,\% \\
Frecuencia de documento (DF) & 0{,}0512 & 17{,}2\,\% \\
Palabra (embedding) & 0{,}0418 & 14{,}1\,\% \\
Token siguiente (POS) & 0{,}0289 & 9{,}7\,\% \\
Es capitalizado & 0{,}0156 & 5{,}3\,\% \\
Lema & 0{,}0089 & 3{,}0\,\% \\
Frecuencia absoluta (TF) & 0{,}0034 & 1{,}1\,\% \\
Longitud de palabra & 0{,}0021 & 0{,}7\,\% \\
Es todo mayúsculas & 0{,}0012 & 0{,}4\,\% \\
\bottomrule
\end{tabular}
\end{table}

\begin{table}[H]
\centering
\small
\caption{Importancia de rasgos en BiLSTM-CRF para español. Mismo formato que Tabla~\ref{tab:bio_importancia_en}.}
\label{tab:bio_importancia_es}
\begin{tabular}{@{} l r r @{}}
\toprule
\textbf{Rasgo} & \textbf{Δ F1 (permutación)} & \textbf{\% Importancia} \\
\midrule
Diversidad contextual (CD) & 0{,}1127 & 31{,}8\,\% \\
Frecuencia de documento (DF) & 0{,}0812 & 22{,}9\,\% \\
Token anterior (POS) & 0{,}0654 & 18{,}5\,\% \\
Palabra (embedding) & 0{,}0521 & 14{,}7\,\% \\
Token siguiente (POS) & 0{,}0326 & 9{,}2\,\% \\
Es capitalizado & 0{,}0142 & 4{,}0\,\% \\
Frecuencia absoluta (TF) & 0{,}0057 & 1{,}6\,\% \\
Lema & 0{,}0041 & 1{,}2\,\% \\
Longitud de palabra & 0{,}0032 & 0{,}9\,\% \\
Es todo minúsculas & 0{,}0018 & 0{,}5\,\% \\
\bottomrule
\end{tabular}
\end{table}

\paragraph{Hallazgo central: CD emerge como predictor primario; TF es marginal.}
Las Tablas~\ref{tab:bio_importancia_en} y~\ref{tab:bio_importancia_es} resumen los 10 rasgos más importantes según BiLSTM-CRF para ambos subcorpus:

El patrón es claro e idéntico en ambas lenguas: \textbf{diversidad contextual (CD)} emerge como el predictor más importante, con contribuciones relativas del 28{,}5\,\% en inglés y 31{,}8\,\% en español. En contraste, \textbf{frecuencia absoluta (TF)}---el número total de ocurrencias--- ocupa posiciones marginales (1{,}1\,\% en EN; 1{,}6\,\% en ES), superada incluso por rasgos morfosintácticos simples (es capitalizado: 5{,}3\,\% EN, 4{,}0\,\% ES).

Con CD operacionalizada como dispersión documental, esto refleja una estructura profunda: una palabra que aparece 1{,}000 veces pero confinada a solo 8 documentos (CD=8) es un predictor mucho más débil que un término que aparece 50 veces pero disperso a través de 45 documentos distintos (CD=45). Operacionalmente, la amplitud de dispersión disciplinaria supera a la abundancia bruta como predictor de terminologicidad.

Este resultado alinea conceptualmente con la literatura: Repar et al. (2025), en ablación controlada sobre RSDO5 esloveno, hallan que las features contextuales (tópicos, patrones) mantienen F1~0{,}35--0{,}49 incluso cuando se remueven features estadísticas (que caen a F1~0), confirmando que rasgos contextuales son más robustos en bajo-recurso/no-visto. La novedad aquí es cuantificar esta robustez explícitamente mediante partición seen/unseen.

\paragraph{Validación estadística: Predicción de términos no vistos por familia de rasgo.}

Para cuantificar la relación entre importancia de rasgos y predicción de términos no vistos, calculamos la correlación de Spearman ($\rho_s$) entre el rank de importancia de cada rasgo (según permutation importance en BiLSTM-CRF) y su poder predictivo sobre términos unseen en el protocolo LOTO (términos que aparecen en test pero nunca en train).

\begin{table}[H]
\centering
\small
\caption{Correlación de Spearman entre importancia de rasgo y detectabilidad de términos no vistos en holdout span-level. Cuantifica si rasgos con mayor importancia según BiLSTM-CRF predicen mejor los términos que el modelo no ha visto durante entrenamiento. Comparación con la literatura: Rigouts Terryn et al. (2021) reportan decay similar (RC contextual ~0.43, RC frecuencia ~−0.15) aunque en test cross-domain.}
\label{tab:correlacion_importancia_detectabilidad}
\begin{tabular}{@{} l r r r @{}}
\toprule
\textbf{Familia de rasgo} & \textbf{EN: $\rho_s$} & \textbf{ES: $\rho_s$} & \textbf{Significancia} \\
\midrule
Diversidad contextual (CD) & 0{,}35 & 0{,}46 & $p < 0{,}001$ ***\\
Frecuencia de documento (DF) & 0{,}18 & 0{,}22 & $p < 0{,}05$ \\
POS contextuales (tokens adj.) & 0{,}12 & 0{,}16 & $p < 0{,}10$ \\
Frecuencia absoluta (TF) & 0{,}02 & 0{,}03 & n.s. ($p = 0{,}14$ EN; $p = 0{,}18$ ES) \\
\bottomrule
\end{tabular}
\end{table}

\begin{figure}[H]
    \centering
    \includegraphics[width=\linewidth]{images/scatter_recall_rasgos.png}
    \caption{Scatter plot: Recall de términos no vistos vs. Diversidad Contextual (CD) por término en ambos subcorpus. Panel izquierdo: inglés; panel derecho: español. La relación positiva con CD es visible y significativa (ρₛ = 0.35 EN, ρₛ = 0.46 ES; p < 0.001), mientras que la relación con TF es plana.}
    \label{fig:scatter_recall_rasgos}
\end{figure}

El resultado es contundente: la correlación entre diversidad contextual y detectabilidad en términos no vistos es altamente significativa ($\rho_s = 0{,}35$ EN, $\rho_s = 0{,}46$ ES; ambos $p < 0{,}001$), mientras que la correlación con frecuencia absoluta es no significativa en ambos subcorpus ($p = 0{,}14$ EN; $p = 0{,}18$ ES). En otras palabras: \textbf{los rasgos contextuales (CD, DF) predicen significativamente mejor los términos nuevos que los rasgos de frecuencia bruta}. Esta validación empírica corrobora directamente lo que Rigouts Terryn et al. observaron indirectamente mediante recall en hapax: rasgos contextuales mitigan el problema de términos raros/no vistos que aquejan a sistemas basados únicamente en frecuencia.

\paragraph{Implicaciones teóricas: Validación empírica de la teoría de Cabré y replicación conceptual de hallazgos en la literatura.}

Estos resultados constituyen una validación empírica sólida de los fundamentos teóricos de la terminología contemporánea, particularmente la concepción de Cabré (1999) del término como unidad contextual antes que léxica: ``los términos no son palabras, sino las palabras cuando funcionan como denominaciones especializadas de conceptos''. La teoría postula que la terminologicidad es una propiedad de \emph{uso discursivo}, no de identidad léxica inherente; una palabra es término en ciertos contextos especializados, pero no término en contextos generales.

Nuestros resultados empíricos dan cuerpo explícitamente a este marco mediante el test LOTO: el modelo aprende a identificar términos no mediante memorización de identidades de palabras frecuentes (donde TF=1{,}1\,\% de importancia), sino mediante patrones de dispersión contextual. Una unidad que aparece dispersa en 40 documentos distintos dentro del dominio especializado (CD=40) adquiere una firma distribucional que el modelo reconoce como característica de términos genuinos. Inversamente, una palabra que aparece 1{,}000 veces pero confinada a solo 3 documentos (CD=3) es débilmente detectada porque carece de amplitud de dispersión disciplinaria.

Este hallazgo **replica conceptualmente** lo observado indirectamente en Rigouts Terryn et al. (2021--2022) y TermEval 2020: que rasgos contextuales mitigan el problema de términos raros (hapax) que colapsa en sistemas puramente frecuenciales. La novedad metodológica es cuantificar esto mediante partición visto/no-visto explícita, que no existe en ATE publicado. Así, el presente trabajo valida la hipótesis teórica de Cabré no solo como intuición sino como predictor cuantificable de generalización a términos nuevos.

Finalmente, esta conclusión resuelve la paradoja de la Zipfian tail (§\ref{subsec:zipfian_tail}): aunque muchos términos aparecen pocas veces (tail de frecuencia), si esas apariciones están dispersas a través de múltiples documentos del dominio, el modelo puede inferir correctamente su terminologicidad sin haber visto esa unidad exacta durante entrenamiento. La clave no es la frecuencia sino la \emph{variedad de contextos}.

\paragraph{Consecuencias para diseño de extractor automático: Rasgos contextuales e integración híbrida.}

El hallazgo tiene implicaciones prácticas directas para sistemas de ATE:

\begin{enumerate}
    \item \textbf{Métricas de contexto son prioritarias:} Cualquier sistema que busque generalización robusta a términos nuevos debe \emph{explícitamente} computar y priorizar métricas de dispersión contextual (CD, document frequency, coocurrence patterns) sobre índices puramente frecuenciales (TF-IDF, C-value). Los índices de terminologicidad que consideran únicamente frecuencia sin medidas de spread contextual están suboptimizados desde el punto de vista del aprendizaje supervisado.

    \item \textbf{Limitaciones de CD y camino hacia embeddings contextuales:} Esta investigación usa CD como feature artesanal simple. La literatura más reciente (TermEval 2020, Rigouts Terryn 2021--2022) demuestra que embeddings contextuales aprendidos (BERT, ELMo, BiLSTM contextuales) superan a features manuales de dispersión. Sin embargo, CD tiene la ventaja de ser interpretable y rápida de computar en aplicaciones de bajo-recurso. Un diseño híbrido robusto combinaría ambas: features de contexto simples (CD, DF) como señal rápida + embeddings contextuales para refinamiento.

    \item \textbf{Trade-off precisión-recall:} Los rasgos contextuales típicamente aumentan el recall en términos no vistos (0{,}43 en hapax aquí, vs 0{,}34 en HAMLET) pero pueden reducir la precisión. Los sistemas de producción deben reportar curvas P-R completas y elegir punto operativo según el caso de uso (máxima precisión vs máximo recall para descubrimiento terminológico).

    \item \textbf{Integración en pipeline híbrido:} Los sistemas más sofisticados (como TECLA en §\ref{sec:tecla}) integran distribucionalismo explícitamente en las fases híbridas de extracción (cálculo de CD como parte de ranking de candidatos) y lo aprovechan downstream en clasificadores supervisados. Esta investigación valida que esta estrategia es empíricamente fundamentada.
\end{enumerate}

\paragraph{Limitaciones del análisis y trabajos futuros.}

Conviene notar que: (i) la mayoría de la evidencia en la literatura sobre rasgos contextuales procede de embeddings aprendidos (BERT), no de métricas simples como CD; (ii) la relación $\rho_s=0{,}35$--0{,}46 entre CD e importancia, aunque significativa, es moderada, sugiriendo que otras variables (longitud de término, ambigüedad, dominio) también influyen; (iii) la efectividad de CD puede no transferir uniformemente a otros dominios o lenguas, especialmente con distribuciones muy distintas de Zipfian. Trabajos futuros deberían (a) comparar CD con embeddings contextuales en el mismo pipeline LOTO, (b) hacer análisis de ablación explicit por familias de rasgos (solo contextuales, solo frecuencia, híbrido) como hace Repar et al. (2025), y (c) validar en dominios y lenguas adicionales.

% ============================================================
%  SUBSECCIÓN: Términos incrustados
% ============================================================

\subsection{Términos incrustados y cadenas terminológicas}
\label{subsec:terminos_incrustados}

Una propiedad estructural de los inventarios terminológicos especializados que el análisis de tipos y frecuencias no revela directamente es la tendencia de ciertas unidades a actuar como \textbf{núcleos léxicos productivos}: términos que, además de funcionar de manera autónoma, aparecen incrustados como componentes de cadenas terminológicas más largas dentro del propio inventario. Esta propiedad ---denominada aquí \textit{incrustación terminológica}--- tiene consecuencias tanto para la evaluación de la representatividad del inventario como para la ponderación estadística de los candidatos, y constituye precisamente el fenómeno que la métrica \textit{C-value} está diseñada para formalizar.

Para el conjunto de términos validados de ambos subcorpus se ha calculado, para cada unidad $t$ del inventario, cuántos términos más largos del mismo inventario la contienen como subsecuencia contigua de tokens (sus \emph{formas contenedoras}), qué fracción de las ocurrencias de $t$ en el corpus corresponde a apariciones autónomas (\emph{standalone}) y qué fracción corresponde a apariciones incrustadas en una cadena más larga (\emph{nested}).

\paragraph{Inventario en inglés.}
De los 134 términos del subcorpus inglés con valores estadísticos completos, 10~unidades (7,5\,\%) presentan al menos una forma contenedora en el inventario; las 124~restantes (92,5\,\%) no participan como componente de ningún otro término validado. Sin embargo, estas 10~unidades concentran el 16,4\,\% del total de ocurrencias terminológicas del subcorpus, por lo que su impacto en los cómputos de frecuencia es desproporcionadamente alto. La Tabla~\ref{tab:terminos_incrustados_en} recoge las unidades con al menos una forma contenedora, ordenadas por número de formas contenedoras.

\begin{table}[H]
\centering
\small
\setlength{\tabcolsep}{4pt}
\renewcommand{\arraystretch}{1.15}
\caption{Términos del subcorpus inglés con al menos una forma contenedora en el inventario validado. \textit{Frec.\ total} = ocurrencias absolutas en el corpus; \textit{Frec.\ autónoma} = ocurrencias en que el término aparece de forma independiente; \textit{\% incrustada} = proporción de ocurrencias incrustadas en una cadena más larga; \textit{FC} = número de formas contenedoras en el inventario.}
\label{tab:terminos_incrustados_en}
\begin{tabular}{@{} l r r r r r @{}}
\toprule
\textbf{Término} & \textbf{Frec.\ total} & \textbf{Frec.\ autónoma} & \textbf{\% incrustada} & \textbf{FC} & \textbf{C-value} \\
\midrule
\textit{corpus}       & 13\,767 &  9\,335 & 32\,\% & 28 & 10\,480{,}0 \\
\textit{annotation}   &  1\,513 &  1\,129 & 25\,\% &  7 &  1\,520{,}1 \\
\textit{frequency}    &  1\,667 &  1\,352 & 19\,\% &  6 &  1\,909{,}9 \\
\textit{type}         &  1\,048 &    782  & 25\,\% &  1 &  1\,978{,}6 \\
\textit{keyword}      &    337  &    198  & 41\,\% &  1 &    651{,}5 \\
\textit{national corpus} & 189 &     75  & 60\,\% &  1 &    194{,}2 \\
\textit{treebank}     &    105  &     64  & 39\,\% &  1 &     40{,}9 \\
\textit{concordance}  &    161  &    145  & 10\,\% &  1 &    207{,}4 \\
\textit{lob}          &    304  &    163  & 46\,\% &  1 &      2{,}4 \\
\textit{parsing}      &     54  &     49  &  9\,\% &  1 &     40{,}1 \\
\bottomrule
\end{tabular}
\end{table}

Los datos ponen de manifiesto patrones de diferente naturaleza. En un extremo, \textit{corpus} actúa como el núcleo más productivo del dominio con diferencia: 28~de los términos validados lo contienen como componente (desde \textit{reference corpus} o \textit{spoken corpus} hasta \textit{British National Corpus} o \textit{diachronic corpus}), y el 32\,\% de sus 13\,767~ocurrencias son incrustadas ---es decir, corresponden a contextos en que el término forma parte de una cadena denominativa más larga---. En un extremo opuesto, \textit{lob} ilustra con especial claridad el efecto de la incrustación sobre la estimación de la terminologicidad: con 304~ocurrencias absolutas y frecuencia de documento de~11, su frecuencia bruta lo situaría entre los candidatos de peso medio del inventario; sin embargo, el 46\,\% de sus apariciones corresponden a la forma contenedora \textit{LOB Corpus}, y su C-value resultante es de~2{,}4, el más bajo entre todos los términos con formas contenedoras. La desconexión entre frecuencia absoluta y C-value en este caso no es un artefacto del sistema, sino la consecuencia directa y correcta del descuento por incrustación: \textit{lob} es terminológicamente poco autónomo porque, en la mayoría de sus contextos de uso, su referencia queda absorbida por una unidad denominativa más precisa.

\paragraph{Inventario en español.}
En el subcorpus español, la incrustación terminológica es considerablemente menos frecuente: solo 5~de los 157~términos validados (3,2\,\%) presentan al menos una forma contenedora, y el conjunto de sus ocurrencias incrustadas representa el 9,0\,\% del total de ocurrencias terminológicas del subcorpus. La Tabla~\ref{tab:terminos_incrustados_es} recoge estas unidades.

\begin{table}[H]
\centering
\small
\setlength{\tabcolsep}{4pt}
\renewcommand{\arraystretch}{1.15}
\caption{Términos del subcorpus español con al menos una forma contenedora en el inventario validado. Mismas definiciones de columna que la Tabla~\ref{tab:terminos_incrustados_en}.}
\label{tab:terminos_incrustados_es}
\begin{tabular}{@{} l r r r r r @{}}
\toprule
\textbf{Término} & \textbf{Frec.\ total} & \textbf{Frec.\ autónoma} & \textbf{\% incrustada} & \textbf{FC} & \textbf{C-value} \\
\midrule
\textit{corpus}               & 11\,749 & 10\,256 & 13\,\% & 26 & 11\,691{,}6 \\
\textit{lingüística de corpus}&    690  &    677  &  2\,\% &  1 &  1\,073{,}0 \\
\textit{monitor}              &     44  &     33  & 25\,\% &  1 &     33{,}0 \\
\textit{Trados}               &     24  &     13  & 46\,\% &  1 &     13{,}0 \\
\textit{Transit}              &      5  &      3  & 40\,\% &  1 &      3{,}0 \\
\bottomrule
\end{tabular}
\end{table}

La asimetría entre ambos subcorpus es lingüísticamente significativa. El inglés forma términos especializados de manera preferente mediante composición sintagmática ---secuencias del tipo \texttt{N\,N} o \texttt{(ADJ|N)*\,N}--- que generan cadenas terminológicas de extensión variable sobre un núcleo nominal compartido; el resultado natural son constelaciones de términos que comparten componentes (todas las variantes de \textit{X corpus} o \textit{corpus X}). El español, en cambio, lexicaliza los mismos conceptos en buena medida mediante derivación morfológica (\textit{lematización}, \textit{tokenización}, \textit{representatividad}) o mediante sintagmas preposicionales (\textit{etiquetado de corpus}, \textit{frecuencia de documento}) que no generan la misma relación de incrustación entre términos del inventario. Así, aunque \textit{corpus} es también el núcleo más productivo del subcorpus español con 26~formas contenedoras, su tasa de incrustación es del 13\,\% ---frente al 32\,\% en inglés---, lo que refleja que la diversificación denominativa en español se apoya menos en la composición nominal y más en la expansión preposicional y la derivación.

\paragraph{Incrustación y C-value.}
La relación entre incrustación terminológica y C-value no es accidental: el C-value \autocite{frantzi1999c} fue diseñado precisamente para abordar este problema. Su fórmula descuenta de la frecuencia bruta de un candidato las ocurrencias que corresponden a su aparición como subcadena de términos más largos, ponderadas por el número de esas formas contenedoras:

\begin{equation}
C\text{-value}(t) =
\begin{cases}
\log_2 |t| \cdot f(t) & \text{si } t \text{ no tiene formas contenedoras,} \\[4pt]
\log_2 |t| \cdot \left( f(t) - \dfrac{1}{|\mathcal{C}(t)|}\displaystyle\sum_{b \in \mathcal{C}(t)} f(b) \right) & \text{en caso contrario,}
\end{cases}
\label{eq:cvalue_incrustacion}
\end{equation}

donde $|t|$ es la longitud en tokens, $f(t)$ la frecuencia bruta, $\mathcal{C}(t)$ el conjunto de formas contenedoras y $f(b)$ la frecuencia de cada una de ellas. El análisis de incrustación presentado en esta sección es, por tanto, la manifestación empírica de lo que el C-value computa internamente: un término con alta frecuencia bruta pero bajo C-value es, casi sin excepción, una unidad que aparece mayoritariamente como componente de términos más largos, como ilustra el caso de \textit{lob} en inglés (frecuencia bruta~304, C-value~2{,}4) o de \textit{Trados} en español (frecuencia bruta~24, C-value~13{,}0, absorbida en un 46\,\% por \textit{SDL Trados Studio}).

En suma, la incrustación terminológica constituye una propiedad discreta pero cuantitativamente relevante de los inventarios especializados: no afecta al mayor número de tipos del vocabulario, pero concentra su efecto en los núcleos conceptuales más productivos del dominio ---precisamente aquellos con mayor capacidad generativa para la formación de cadenas denominativas nuevas---. Su formalización mediante el C-value permite que la puntuación de terminologicidad final (§\ref{sec:extraccion}) refleje la autonomía referencial de cada unidad y no solo su peso distribucional bruto.

% ── Modelo formal del término (OE1.3) ──────────────────────────────────────

\paragraph{El modelo teórico-computacional del término.}
Los resultados de la extracción permiten articular de forma explícita el \textbf{modelo teórico-computacional del término} al que apunta el objetivo específico~OE1.3: una representación formal que integra los rasgos identificados en una abstracción operacionalizable. En el marco de esta investigación, una unidad léxica $u$ se caracteriza como término a través de tres dimensiones complementarias:

\begin{enumerate}
    \item \textbf{Dimensión formal — \textit{unithood}.} $u$ satisface al menos uno de los patrones morfosintácticos de la gramática terminológica $\mathcal{G}$ (Tabla~\ref{tab:patrones-pos}), garantizando que posee la estructura sintagmática adecuada para actuar como unidad denominativa:
    \begin{equation}
        \mathrm{unithood}(u) = \mathbf{1}\!\left[\exists\, p \in \mathcal{G} : u \models p\right]
        \label{eq:unithood}
    \end{equation}

    \item \textbf{Dimensión estadística — \textit{termhood}.} La terminologicidad de $u$ en el corpus se cuantifica mediante el índice compuesto (ec.~\ref{eq:terminologicidad}), que combina el peso distribucional TF-IDF con la autonomía sintagmática capturada por el C-value:
    \begin{equation}
        T(u) = 0{,}6\cdot\widetilde{\mathrm{TF\text{-}IDF}}(u) + 0{,}4\cdot\widetilde{C\text{-value}}(u)
        \label{eq:termhood}
    \end{equation}
    siendo $T(u) > \theta_{90}$ ---el percentil~90 de la distribución--- la condición de umbral para la validación experta.

    \item \textbf{Dimensión distribucional — gradiente de terminologicidad.} Los rasgos distribucionales del lema de $u$ en el corpus ---en especial su diversidad contextual ($\mathrm{CD}$), su categoría gramatical y su frecuencia de documento--- conforman un vector de propiedades que el clasificador supervisado aprende a asociar con la etiqueta de término independientemente de la identidad léxica. El análisis de importancia (Tablas~\ref{tab:bio_importancia_en} y~\ref{tab:bio_importancia_es}) muestra que $\mathrm{CD}(u)$ es el predictor primario en ambas lenguas, resultado que se confirma también en el holdout ($\rho_s \geq 0{,}35$, $p < 0{,}001$): los términos con mayor diversidad de contextos discursivos son sistemáticamente más detectables como unidades terminológicas incluso cuando nunca han sido etiquetados durante el entrenamiento.
\end{enumerate}

La combinación de las tres dimensiones define una caracterización del término como unidad que (i)~cumple restricciones formales estructurales, (ii)~posee un peso distribucional estadísticamente marcado en el corpus y (iii)~aparece en una variedad suficientemente heterogénea de contextos discursivos. Esta triple condición articula computacionalmente la concepción terminológica contemporánea del término como unidad \emph{formal}, \emph{funcional} y \emph{contextual} \cite{cabre1999terminology}, y constituye el modelo teórico-computacional que sustenta tanto la herramienta TECLA como los experimentos de extracción supervisada del capítulo.

La sección siguiente traslada el foco desde la caracterización del inventario terminológico hacia su detección automática: ¿en qué medida pueden los clasificadores supervisados aprender, a partir de rasgos lingüísticos y estadísticos, qué propiedades hacen que una unidad léxica sea terminológica, y cuánto de ese aprendizaje es genuinamente transferible a términos que el modelo nunca ha visto etiquetados?

\section{TECLA: \emph{Term Extraction and Corpus Linguistics Application}}
\label{sec:tecla}

La librería TECLA (\textit{Term Extraction and Corpus Linguistics Application}) se articula como una estación de trabajo integral para lingüistas y especialistas en terminología que deseen abordar de forma sistemática el análisis terminológico a partir de la explotación de corpus. Su diseño responde a la necesidad de contar con un entorno unificado que reúna en un único sistema todas las fases del trabajo terminológico automatizado: desde la compilación y descripción del corpus hasta la extracción de candidatos, la evaluación de la representatividad terminológica y la detección supervisada de unidades terminológicas. El objetivo es proporcionar una herramienta flexible, transparente y accesible que combine metodologías lingüísticas y computacionales sin renunciar a los principios de interpretabilidad, trazabilidad y control por parte del investigador. Frente a estrategias fragmentadas o centradas exclusivamente en criterios estadísticos, TECLA integra tanto enfoques basados en reglas lingüísticas como métricas distribucionales, permitiendo articular flujos completos de trabajo terminológico con salida replicable y personalizable.

\subsection{Arquitectura}
\label{subsec:arquitectura_tecla}

La arquitectura de TECLA se organiza en cuatro módulos funcionales interoperables, que pueden emplearse de forma independiente o encadenada dentro de un flujo de trabajo completo. A continuación se describe cada módulo con el nivel de detalle necesario para evaluar tanto sus capacidades como sus limitaciones metodológicas.

\subsubsection*{Módulo 1. Gestión del corpus}

Este módulo cubre las operaciones previas a la extracción terminológica propiamente dicha: la constitución, limpieza, enriquecimiento y descripción del corpus de trabajo. Acepta como entrada documentos en formato \texttt{PDF} o texto plano (\texttt{.txt}); en el primer caso, realiza la conversión a texto plano mediante \texttt{pdfminer.six} y deposita los ficheros resultantes en un directorio de trabajo estructurado (\texttt{\_tecla\_txt/raw/} para el texto sin procesar, \texttt{\_tecla\_txt/clean/} para el texto limpiado).

La \textbf{detección de lengua} puede configurarse manualmente (\texttt{es}/\texttt{en}) o delegarse al módulo de detección automática basado en \texttt{langid}, que identifica la lengua dominante de cada documento y selecciona el modelo spaCy correspondiente (\texttt{es\_core\_news\_md} para el español; \texttt{en\_core\_web\_lg} para el inglés).

El \textbf{módulo de limpieza} implementa más de una docena de operaciones configurables de forma independiente mediante banderas booleanas, lo que garantiza que nada se aplica implícitamente y que el usuario tiene control total sobre cada transformación:

\begin{itemize}
    \item \textit{Ruido estructural}: reunión de palabras partidas por guión al final de línea; eliminación de números de página y líneas muy cortas.
    \item \textit{Metadatos documentales}: eliminación de títulos de sección aislados (\textit{Abstract}, \textit{Introducción}, \textit{Methodology}…), líneas de palabras clave, pies de figura y tabla, líneas de URL/DOI/email, marcadores de copyright y licencias, datos de publicación de revista (volumen, número, páginas), marcadores de cita numérica en línea (\texttt{[1]}, \texttt{[2,3]}) y líneas mayoritariamente numéricas (tablas estadísticas).
    \item \textit{Sección de referencias}: eliminación automática del bloque bibliográfico al final del documento.
    \item \textit{Filtrado léxico}: eliminación de palabras vacías (combinando las listas de spaCy y NLTK con una stoplist adicional personalizable por el usuario) y lematización opcional mediante spaCy.
\end{itemize}

La \textbf{inferencia y edición de metadatos} permite asociar a cada documento información bibliográfica (autor, año, tema, fuente, notas libres) tanto por inferencia automática a partir del nombre de archivo como mediante edición manual en la interfaz gráfica. Los metadatos se almacenan en un fichero \texttt{metadata.csv} vinculado al directorio de trabajo.

Las \textbf{estadísticas descriptivas} se calculan tanto a nivel de documento como de corpus global e incluyen: número de tokens, número de tipos, relación tipo-token (TTR), proporción de hápax, densidad léxica, número de oraciones y longitud media de oración. Todos estos indicadores pueden exportarse a CSV para su análisis externo.

\subsubsection*{Módulo 2. Extracción de candidatos a término}

Este módulo implementa el núcleo metodológico del sistema de extracción terminológica no supervisada. Sigue una arquitectura híbrida en dos etapas secuenciales:

\paragraph{Filtro lingüístico.} La primera etapa aplica un conjunto de patrones morfosintácticos definidos en un fichero de gramática (\texttt{grammar.yaml}) sobre los documentos procesados por spaCy. Cada patrón especifica una secuencia de etiquetas Universal Dependencies (UD) que debe satisfacer un candidato para ser considerado estructuralmente plausible como término (\textit{unithood}). Los patrones son configurables: el sistema incluye una gramática de referencia para inglés y español, pero el usuario puede substituirla por una gramática propia adaptada a su dominio. La Tabla~\ref{tab:patrones_extraccion} recoge los patrones incorporados, con su descripción estructural y ejemplos.

\begin{table}[H]
\centering
\small
\begin{tabular}{@{} l l p{4.5cm} p{4cm} @{}}
\toprule
\textbf{ID} & \textbf{L} & \textbf{Patrón UD} & \textbf{Ejemplos} \\
\midrule
\texttt{EN\_NP\_PRENOMINAL}  & EN & \texttt{(ADJ|NOUN|PROPN)* NOUN}   & \textit{balanced corpus}, \textit{written corpus} \\
\texttt{EN\_NP\_OF}          & EN & \texttt{NOUN+ of NOUN+}           & \textit{part of speech}, \textit{analysis of variance} \\
\texttt{EN\_NP\_PROPN}       & EN & \texttt{PROPN PROPN+}             & \textit{Brown Corpus}, \textit{WordNet 3.0} \\
\texttt{EN\_NOUN\_SEQUENCE}  & EN & \texttt{NOUN NOUN+}               & \textit{token frequency}, \textit{corpus size} \\
\midrule
\texttt{ES\_NP\_POSTNOMINAL} & ES & \texttt{NOUN ADJ+}                & \textit{corpus paralelo}, \textit{frecuencia léxica} \\
\texttt{ES\_NP\_DE}          & ES & \texttt{NOUN+ de NOUN+}           & \textit{etiquetado de corpus}, \textit{lista de palabras} \\
\texttt{ES\_NP\_DE\_ADJ}     & ES & \texttt{NOUN+ de NOUN ADJ+}       & \textit{análisis de varianza estadística} \\
\texttt{ES\_NP\_PRENOMINAL}  & ES & \texttt{(ADJ|NOUN|PROPN)* NOUN}   & \textit{nueva lingüística computacional} \\
\texttt{ES\_NP\_PROPN}       & ES & \texttt{PROPN PROPN+}             & \textit{corpus CREA}, \textit{BNC} \\
\texttt{ES\_NP\_NOUN\_PROPN} & ES & \texttt{NOUN PROPN+}              & \textit{corpus CORDE} \\
\bottomrule
\end{tabular}
\caption{Patrones morfosintácticos incorporados en la gramática de TECLA para inglés (EN) y español (ES). Las etiquetas POS siguen el estándar Universal Dependencies. El operador \texttt{*} indica cero o más ocurrencias; \texttt{+} indica una o más.}
\label{tab:patrones_extraccion}
\end{table}

Adicionalmente, el filtro aplica dos mecanismos de depuración: un \textbf{filtro de palabras funcionales} que rechaza candidatos que contienen determinantes, pronombres, auxiliares, conjunciones o interjecciones en cualquier posición; y un \textbf{filtro de artefactos} que elimina segmentos producidos por errores de tokenización (palabras partidas por guión que no se reunificaron) y candidatos de alta frecuencia que aparecen en un único documento (ruido de índices o metadatos repetidos).

\paragraph{Filtro estadístico.} La segunda etapa enriquece la lista de candidatos ya depurada con métricas de terminologicidad (\textit{termhood}). Para cada candidato se calculan:

\begin{itemize}
    \item \textbf{Frecuencia absoluta} ($f$), \textbf{frecuencia de documento} (número de documentos en que aparece) y \textbf{frecuencia relativa} (ocurrencias por mil tokens).
    \item \textbf{TF-IDF} (frecuencia de término multiplicada por el logaritmo inverso de la frecuencia de documento): mide hasta qué punto una unidad es característica del corpus frente a ser omnipresente.
    \item \textbf{C-value} \autocite{frantzi1999c}: métrica diseñada específicamente para términos multipalabra que penaliza a los candidatos cortos que aparecen frecuentemente como subcadenas de candidatos más largos; formaliza la noción de que un candidato frecuente incrustado en muchos términos más largos tiene menos autonomía terminológica.
    \item \textbf{NC-value} \autocite{frantzi2000automatic}: extensión contextual del C-value que incorpora la información proporcionada por las palabras de contexto (±2 tokens) de cada ocurrencia del candidato. Cuanto más frecuentemente aparece una palabra de contenido como contexto de múltiples candidatos, mayor peso recibe, lo que permite discriminar candidatos terminológicos en función de su entorno discursivo habitual.
\end{itemize}

Los tres índices numéricos (TF-IDF, C-value, NC-value) se normalizan independientemente al intervalo $[0,1]$ mediante normalización min-max. La puntuación de \textbf{terminologicidad} final es un promedio ponderado de TF-IDF$_{\text{norm}}$ y NC-value$_{\text{norm}}$:

\begin{equation}
\text{terminologicidad} = \alpha \cdot \text{TF-IDF}_{\text{norm}} + (1-\alpha) \cdot \text{NC-value}_{\text{norm}}
\label{eq:terminologicidad_tecla}
\end{equation}

El parámetro $\alpha$ determina el perfil de puntuación y puede ajustarse mediante tres presets predefinidos: \textit{Frecuencia} ($\alpha=0{,}6$), orientado a términos con alta presencia en el corpus; \textit{Equilibrado} ($\alpha=0{,}5$), que pondera por igual ambas métricas; y \textit{Complejidad} ($\alpha=0{,}2$), que prioriza los términos multipalabra con estructuras complejas y bien delimitadas por el contexto.

\paragraph{Postprocesamiento.} Antes de presentar el inventario final al usuario, el sistema aplica tres operaciones adicionales:

\begin{itemize}
    \item \textbf{Detección de variantes}: las variantes ortográficas (diferencias de mayúsculas, guionización o espaciado, como \textit{machine learning} / \textit{machine-learning}) se agrupan bajo la forma canónica con mayor terminologicidad; las siglas cuyas iniciales coinciden con las palabras de un término más largo se vinculan bidireccionalmente (p.\,ej., \textit{NLP} ↔ \textit{natural language processing}).
    \item \textbf{Realce de núcleos}: los unigramas que actúan como núcleo sintáctico de varios términos multipalabra reciben un incremento proporcional de su puntuación de terminologicidad (hasta un máximo del 35\,\%), con el fin de compensar la pérdida de señal de los sustantivos nucleares de alta productividad terminológica.
    \item \textbf{Filtrado de términos genéricos}: una lista curada de expresiones académicas y metodológicas que superan los filtros morfosintácticos pero no denotan conceptos especializados (\textit{research method}, \textit{resultados obtenidos}, etc.) se elimina antes de presentar los resultados.
\end{itemize}

Los parámetros de extracción configurables por el usuario incluyen: longitud máxima del candidato (1--15 tokens), frecuencia mínima (1--50), umbral de percentil (\textit{p}25, \textit{p}50, \textit{p}75, \textit{p}90) y umbral de score mínimo. La salida puede exportarse en formato \textbf{CSV} (con todos los campos numéricos: término, lema, variantes, patrón, frecuencias, TF-IDF, C-value, NC-value, terminologicidad) o en formato \textbf{TBX-Basic} (ISO 30042), el estándar de intercambio terminológico utilizado por herramientas como SDL MultiTerm, OmegaT y Terminator, lo que facilita la integración del inventario en flujos de trabajo terminográficos profesionales.

\subsubsection*{Módulo 3. Evaluación de la representatividad terminológica}

Este módulo permite al investigador responder a una pregunta previa a cualquier tarea terminológica: \emph{¿es este corpus suficientemente representativo del dominio para extraer de él un inventario terminológico fiable?} Para ello implementa el Índice de Representatividad Terminológica (ITR), ya descrito en detalle en el Capítulo~\ref{chap:3_extraccion}, que sintetiza cuatro métricas complementarias en una puntuación única en la escala 0--100:

\begin{itemize}
    \item \textbf{ITC\textsubscript{norm}} (curva de cobertura incremental terminológica): mide la velocidad con que el corpus acumula nuevos tipos de término al añadir documentos sucesivos; un corpus saturado muestra una curva que se aplana rápidamente.
    \item \textbf{MATTR} (\textit{moving average type-token ratio}): estima la riqueza y variedad léxica del corpus mediante una ventana deslizante que evita el efecto de longitud del TTR clásico.
    \item \textbf{Exponente de Heaps} ($\beta$): parámetro de la ley de Heaps que describe la velocidad de crecimiento del vocabulario con el tamaño del corpus; valores bajos de $\beta$ indican alta densidad terminológica.
    \item \textbf{CMN} (cobertura media normalizada): estimación bootstrapped de qué proporción del inventario terminológico total del dominio cubre ya el corpus.
\end{itemize}

La escala de interpretación del ITR es la siguiente: $\geq80$ (muy representativo), $65$--$79$ (representativo), $50$--$64$ (moderadamente representativo), $35$--$49$ (poco representativo) y $<35$ (escasamente representativo).

Además del ITR global, el módulo proporciona dos análisis complementarios de particular interés metodológico:

\begin{itemize}
    \item \textbf{Dispersión de Gries (DP)}: para cada candidato a término se calcula su coeficiente de dispersión entre los documentos del corpus \autocite{gries2021dispersion}. Un DP próximo a 0 indica que el término aparece de forma uniforme en todos los documentos (señal distribucional robusta); un DP próximo a 1 indica concentración en uno o pocos documentos, lo que puede apuntar a un término de autor o a ruido. Los 50 candidatos con mayor terminologicidad se muestran en una tabla ordenada por DP, lo que permite identificar candidatos problemáticos antes de la validación manual.
    \item \textbf{Keyness G\textsuperscript{2}}: análisis de contraste léxico de los candidatos del corpus analizado frente a un corpus de referencia externo, cargado opcionalmente por el usuario. Para cada candidato se calcula el estadístico $G^2$ de log-verosimilitud \autocite{dunning1993accurate} y el log-ratio, lo que permite identificar qué unidades son genuinamente características del dominio especializado y cuáles podrían pertenecer al léxico general o a otro subdominio.
\end{itemize}

El módulo permite asimismo cargar una \textbf{lista de referencia terminológica} externa (p.\,ej., un glosario publicado, una base de datos de terminología como IATE, o el propio gold standard del investigador) y calcular qué proporción de esa lista aparece entre los candidatos extraídos, produciendo automáticamente métricas de precisión y recall respecto al estándar externo.

\subsubsection*{Módulo 4. Clasificación supervisada}

Este módulo permite entrenar clasificadores automáticos de detección terminológica a partir de un corpus anotado mediante la codificación BIO. El flujo metodológico incluye la extracción de rasgos lingüísticos y distribucionales para cada token del corpus, la organización de los datos de entrenamiento y evaluación mediante particiones cruzadas por documento, y la aplicación de algoritmos de aprendizaje supervisado. Los detalles del espacio de rasgos, los algoritmos implementados, la estrategia de evaluación y los resultados obtenidos en el corpus del presente trabajo se describen en detalle en la sección anterior (\S\ref{sec:modelo_tc}). En el contexto del flujo de trabajo de TECLA, este módulo ocupa la posición final de la cadena: el \textit{Silver Standard} producido por el módulo de extracción sirve como punto de partida para la anotación supervisada, y el clasificador resultante puede aplicarse como un filtro adicional de alta precisión sobre nuevos corpus del mismo dominio.

\paragraph{Disponibilidad en el release.} El módulo de clasificación supervisada, incluyendo la implementación de redes CRF y BiLSTM-CRF, forma parte integral del código fuente de TECLA y fundamenta los resultados de evaluación del presente trabajo. Sin embargo, su distribución en el release público se ha circunscrito al código fuente documentado sin incluir los pesos preentrenados de los modelos. Esta decisión responde a consideraciones prácticas: (i) los modelos neuronales BiLSTM entrenan sobre GPU y requieren dependencias adicionales (PyTorch, CRFsuite compilado) que aumentan significativamente el peso de instalación y la complejidad de configuración; (ii) el público objetivo de la herramienta (investigadores en lingüística, terminólogos, lingüistas computacionales sin experiencia en deep learning) puede extraer terminología con alta calidad mediante los módulos 1--3 sin necesidad de entrenar clasificadores; (iii) la calibración de estos modelos es dominio-dependiente y requiere anotación manual, por lo que un release con modelos preentrenados sería engañoso en su generalización a nuevos domios. El release estable de TECLA se enfoca, por tanto, en consolidar los módulos de extracción no supervisada (módulo 2), evaluación de representatividad (módulo 3) e interfaz de usuario, que ofrecen la máxima relación entre utilidad y accesibilidad. Los investigadores interesados en reproducir los experimentos de clasificación supervisada pueden consultar la documentación del repositorio para instrucciones de instalación de dependencias opcionales y reproducción del pipeline completo.

\subsubsection*{Interfaz gráfica y línea de comandos}

TECLA ofrece dos modos de interacción: una interfaz gráfica de escritorio y una interfaz de línea de comandos.

La \textbf{interfaz gráfica} (GUI) está construida con \texttt{customtkinter} y organiza el flujo de trabajo en cuatro pestañas:

\begin{enumerate}
    \item \textbf{Corpus}: muestra una tabla con todos los documentos del corpus (identificador, tokens, oraciones, tipos, TTR, autor, año, tema) y un panel de edición de metadatos por documento. Incluye exportación de estadísticas a CSV.
    \item \textbf{Extracción}: registra el proceso de extracción en tiempo real en un log y presenta la tabla de candidatos a término ordenable por cualquier métrica (score, TF-IDF, C-value, NC-value, frecuencia absoluta, frecuencia relativa, número de documentos). Las columnas visibles son configurables. Permite exportar la selección a CSV o TBX-Basic.
    \item \textbf{Evaluación}: muestra el score ITR con escala visual interpretada, la curva de crecimiento terminológico, la tabla de dispersión DP para los 50 candidatos más prominentes y el análisis de keyness G\textsuperscript{2} frente a un corpus de referencia externo (opcional). Calcula y muestra la cobertura de la extracción respecto a una lista de referencia terminológica externa (opcional).
    \item \textbf{Revisión}: módulo de validación manual que permite al investigador o terminógrafo clasificar cada candidato como \textit{aprobado}, \textit{rechazado} o \textit{pendiente}, añadir notas libres por término, filtrar la lista por estado de revisión y consultar las estadísticas completas de cada candidato (score, métricas individuales, variantes, patrón morfosintáctico). Esta pestaña implementa la interfaz de trabajo propia del terminógrafo, integrando la labor experta en el mismo entorno que el resto del análisis.
\end{enumerate}

La sesión se persiste automáticamente entre ejecuciones (fichero \texttt{.tecla\_session.json}), de modo que el investigador puede retomar el trabajo de extracción y revisión exactamente en el punto en que lo dejó.

La \textbf{interfaz de línea de comandos} (CLI), invocable mediante \texttt{tecla <subcomando>}, expone tres operaciones principales: \texttt{extract} (extracción completa con opciones de filtrado, exportación CSV y TBX), \texttt{stats} (estadísticas descriptivas del corpus) y \texttt{evaluate} (comparación de la extracción frente a un gold standard con cálculo de precisión, recall y F1). La CLI permite integrar TECLA en flujos de trabajo automatizados o en entornos sin interfaz gráfica.

La Tabla~\ref{tab:dependencias_tecla} recoge las dependencias del paquete.

\begin{table}[H]
\centering
\fontsize{10}{12}\selectfont
\begin{tabular}{@{}p{3.5cm}p{8cm}p{2cm}@{}}
\toprule
\textbf{Librería} & \textbf{Funcionalidad principal} & \textbf{Versión} \\
\midrule
\texttt{spaCy} & Tokenización, lematización y etiquetado morfosintáctico (POS) & $\geq$3.0 \\
\texttt{pdfminer.six} & Extracción de texto a partir de documentos PDF & -- \\
\texttt{langid} & Detección automática de la lengua del corpus & -- \\
\texttt{nltk} & Listas de palabras vacías complementarias & -- \\
\texttt{pyyaml} & Carga de la gramática POS desde fichero YAML & -- \\
\texttt{pandas} & Manipulación y exportación de datos tabulares & $\geq$2.0 \\
\texttt{numpy} & Operaciones numéricas y normalización de scores & $\geq$1.24 \\
\texttt{matplotlib} & Visualización de curvas de crecimiento y métricas & $\geq$3.7 \\
\texttt{tqdm} & Barra de progreso para el procesamiento de textos & $\geq$4.65 \\
\texttt{customtkinter} & Interfaz gráfica de escritorio (GUI) & -- \\
\bottomrule
\end{tabular}
\caption{Dependencias del paquete TECLA. Todas las bibliotecas son de código abierto e instalables mediante \texttt{pip}. Las bibliotecas sin versión mínima especificada son compatibles con cualquier versión estable reciente.}
\label{tab:dependencias_tecla}
\end{table}

\subsection{Evaluación}
\label{subsec:evaluacion_tecla}

La evaluación de TECLA se articula en dos planos complementarios: la verificación de componentes y la validación sobre datos reales del dominio.

A nivel de componente, la librería cuenta con un conjunto de pruebas unitarias que cubren las principales transformaciones de datos y cálculos numéricos de cada módulo: correcta conversión y limpieza de documentos, reproducibilidad de los índices estadísticos (TF-IDF, C-value, ITR), integridad del etiquetado BIO y estabilidad de los resultados de clasificación ante distintos entornos de ejecución. Estas pruebas garantizan que el comportamiento de la librería es predecible y replicable antes de su aplicación sobre cualquier corpus.

A nivel de sistema, la evaluación principal de TECLA coincide con el experimento de clasificación supervisada descrito en la sección anterior~(\S\ref{sec:modelo_tc}): los modelos de detección terminológica implementados en el módulo de clasificación automática se evalúan sobre los subcorpus EN y ES mediante validación cruzada con agrupación por documento. Los resultados confirman la adecuación del espacio de representación definido por TECLA: la regresión logística alcanza un F1 de término de 0,83 en inglés y 0,87 en español, y el BiLSTM-CRF supera el 0,96 en ambas lenguas. El protocolo de reserva de tipos terminológicos (\S\ref{tab:holdout_en}) añade una dimensión crítica a esta evaluación: demuestra que los clasificadores entrenados sin rasgos léxicos directos retienen cierta capacidad de generalización a términos nunca vistos, con un recall de entre 0,37 y 0,50 en inglés y entre 0,37 y 0,45 en español.

Por último, la librería ha sido sometida a una evaluación de usabilidad orientada al perfil de usuario objetivo ---investigadores en lingüística y terminología con conocimientos básicos de Python---. Los resultados confirmaron que la documentación generada con Sphinx y la \textsc{api} de la librería permiten reproducir los experimentos y adaptar los parámetros a corpus propios sin conocimientos especializados en aprendizaje automático.

\section{Glosario digital}
\label{sec:glosario}

Como resultado del proceso de extracción y validación descrito a lo largo de este capítulo, se ha elaborado un glosario terminológico bilingüe del dominio de la lingüística de corpus, que constituye uno de los productos finales del presente trabajo. El glosario reúne 155 tipos terminológicos en inglés y 146 en español, validados manualmente a partir de la lista de candidatos generada por TECLA y contrastados con fuentes de referencia especializadas.

Los términos se organizan en siete categorías temáticas que articulan los principales ámbitos conceptuales de la disciplina: \textit{Fundamentos y disciplinas} (conceptos constitutivos de la lingüística de corpus como campo), \textit{Diseño y compilación de corpus} (terminología asociada a la construcción y tipología de corpus), \textit{Recursos y corpus de referencia} (denominaciones de corpus y recursos de uso generalizado), \textit{Procesamiento y anotación} (terminología del análisis lingüístico computacional), \textit{Estadística, léxico y búsqueda} (métricas, índices y herramientas de consulta), \textit{Herramientas y software} (aplicaciones de análisis de corpus) y \textit{Tecnología y formatos} (estándares de codificación y formatos de intercambio).

El glosario incluye pares de equivalencia bilingüe cuando el término dispone de denominación consolidada en ambas lenguas: por ejemplo, \textit{corpus linguistics} / \textit{lingüística de corpus}, \textit{lemmatisation} / \textit{lematización} o \textit{concordance} / \textit{concordancia}. Cuando un concepto carece de denominación estabilizada en una de las dos lenguas, se registra como laguna terminológica. El glosario identifica 118 lagunas en el español ---términos del inglés sin equivalente consolidado en español, como \textit{corpus-based}, \textit{corpus-driven} o \textit{corpus size}--- y 97 lagunas en el inglés ---términos del español sin equivalente acuñado, como determinadas acepciones de \textit{corpus telemático} o \textit{corpus oportunista}---. Estas asimetrías reflejan el desequilibrio histórico en la producción bibliográfica de la disciplina: la mayor parte de los textos fundacionales y la terminología técnica de referencia fueron acuñados en inglés, de modo que numerosos conceptos metodológicos carecen de traducción estandarizada al español, lo que constituye en sí mismo un dato de interés para el estudio de la dinámica terminológica del campo.

El glosario digital se pone a disposición de la comunidad investigadora como recurso abierto en formato estructurado (\texttt{.csv}), e incluye para cada entrada la forma canónica del término, su categoría temática, los equivalentes bilingües cuando existen y la indicación de laguna cuando procede.
